%======================================================================================
% !!! GENERATED FILE -- DO NOT EDIT !!!
%
% Generated by scripts/extract_metrics.py from paper/metrics_map.yaml
% Every value below was extracted mechanically from a results file. To change a number,
% change the experiment or the map -- never this file. Edits here are silently destroyed
% on the next run and, worse, break the chain of custody the provenance lines promise.
%
% Regenerate:
%   python scripts/extract_metrics.py --map paper/metrics_map.yaml \
%       --results-root results --out <this file>
% Verify (exits 1 if the results moved and the paper did not):
%   python scripts/extract_metrics.py --map paper/metrics_map.yaml \
%       --results-root results --out <this file> --check
%
% Usage notes:
%   * mean_std and ge macros expand to math-mode material (\pm, \ge): write $\valFoo$.
%   * range macros expand to '1.2--3.4' and belong in text mode.
%   * The provenance line above each macro is the audit trail: file, filter, column,
%     full unrounded value, and n. Rounding here loses nothing an auditor needs.
%   * No timestamp is emitted on purpose: it would churn git and defeat --check.
% 282 macros from 15 results file(s).
%======================================================================================

% note: *** GUARDED — see claim-guards.md G2. THIS MACRO MUST NOT BE A HEADLINE. *** It exists so that if the number appears at all (e.g. an appendix table of raw metrics) it is traceable, NOT so that it can be quoted as "3951 ± 3561 mV". Per-seed: 4122.7 / 229.1 / 1134.0 / 5056.6 / 9213.4 — a 40x spread; the distribution is nowhere near normal and the std exceeds any meaningful voltage scale. This is a diverged integration, and rmse_safe (src/metrics.jl:20-25) drops non-finite pairs, so it is the RMSE of the *finite part* of a blown-up solve, over an unreported subset of timesteps. The honest headline for NODE failure is NodeForecastRollout. Prose: "diverges (V RMSE 10^2–10^4 mV across seeds)".
% source=results/metrics_baseline_multiseed.csv filter={model=NODE, observed=full, window=forecast, metric=V_rmse} column=mean/std raw={mean=3951.178018885416, std=3561.340803031814} n=5
\newcommand{\valNodeForecastVRmse}{3951 \pm 3561}

% note: *** GUARDED — see claim-guards.md G2. THIS MACRO MUST NOT BE A HEADLINE. *** It exists so that if the number appears at all (e.g. an appendix table of raw metrics) it is traceable, NOT so that it can be quoted as "3951 ± 3561 mV". Per-seed: 4122.7 / 229.1 / 1134.0 / 5056.6 / 9213.4 — a 40x spread; the distribution is nowhere near normal and the std exceeds any meaningful voltage scale. This is a diverged integration, and rmse_safe (src/metrics.jl:20-25) drops non-finite pairs, so it is the RMSE of the *finite part* of a blown-up solve, over an unreported subset of timesteps. The honest headline for NODE failure is NodeForecastRollout. Prose: "diverges (V RMSE 10^2–10^4 mV across seeds)".
% source=results/metrics_baseline_multiseed.csv filter={model=NODE, observed=full, window=forecast, metric=V_rmse} column=mean/std raw={mean=3951.178018885416, std=3561.340803031814} n=5
\newcommand{\valNodeForecastVRmseMeanStd}{3951 \pm 3561}

% note: *** GUARDED — see claim-guards.md G2. THIS MACRO MUST NOT BE A HEADLINE. *** It exists so that if the number appears at all (e.g. an appendix table of raw metrics) it is traceable, NOT so that it can be quoted as "3951 ± 3561 mV". Per-seed: 4122.7 / 229.1 / 1134.0 / 5056.6 / 9213.4 — a 40x spread; the distribution is nowhere near normal and the std exceeds any meaningful voltage scale. This is a diverged integration, and rmse_safe (src/metrics.jl:20-25) drops non-finite pairs, so it is the RMSE of the *finite part* of a blown-up solve, over an unreported subset of timesteps. The honest headline for NODE failure is NodeForecastRollout. Prose: "diverges (V RMSE 10^2–10^4 mV across seeds)".
% source=results/metrics_baseline_multiseed.csv filter={model=NODE, observed=full, window=forecast, metric=V_rmse} column=mean raw=3951.178018885416 n=5
\newcommand{\valNodeForecastVRmseMean}{3951}

% note: *** GUARDED — see claim-guards.md G2. THIS MACRO MUST NOT BE A HEADLINE. *** It exists so that if the number appears at all (e.g. an appendix table of raw metrics) it is traceable, NOT so that it can be quoted as "3951 ± 3561 mV". Per-seed: 4122.7 / 229.1 / 1134.0 / 5056.6 / 9213.4 — a 40x spread; the distribution is nowhere near normal and the std exceeds any meaningful voltage scale. This is a diverged integration, and rmse_safe (src/metrics.jl:20-25) drops non-finite pairs, so it is the RMSE of the *finite part* of a blown-up solve, over an unreported subset of timesteps. The honest headline for NODE failure is NodeForecastRollout. Prose: "diverges (V RMSE 10^2–10^4 mV across seeds)".
% source=results/metrics_baseline_multiseed.csv filter={model=NODE, observed=full, window=forecast, metric=V_rmse} column=std raw=3561.340803031814 n=5
\newcommand{\valNodeForecastVRmseStd}{3561}

% note: *** GUARDED — see claim-guards.md G2. THIS MACRO MUST NOT BE A HEADLINE. *** It exists so that if the number appears at all (e.g. an appendix table of raw metrics) it is traceable, NOT so that it can be quoted as "3951 ± 3561 mV". Per-seed: 4122.7 / 229.1 / 1134.0 / 5056.6 / 9213.4 — a 40x spread; the distribution is nowhere near normal and the std exceeds any meaningful voltage scale. This is a diverged integration, and rmse_safe (src/metrics.jl:20-25) drops non-finite pairs, so it is the RMSE of the *finite part* of a blown-up solve, over an unreported subset of timesteps. The honest headline for NODE failure is NodeForecastRollout. Prose: "diverges (V RMSE 10^2–10^4 mV across seeds)".
% source=results/metrics_baseline_multiseed.csv filter={model=NODE, observed=full, window=forecast, metric=V_rmse} column=n_seeds raw=5 n=5
\newcommand{\valNodeForecastVRmseN}{5}

% note: THE headline UDE forecast accuracy. Full-state observation, 358 forecast points (t > 30 ms) of one uninterrupted rollout from t=0 — the model is never re-initialized at t_train_end, so this is extrapolation error, not one-step error. Scored against the CLEAN trajectory. Pair with UdeTrainVRmseVsNoisy for the noise floor.
% source=results/metrics_baseline_multiseed.csv filter={model=UDE, observed=full, window=forecast, metric=V_rmse} column=mean/std raw={mean=0.24523933188134714, std=0.028375840021892113} n=5
\newcommand{\valUdeForecastVRmse}{0.245 \pm 0.028}

% note: THE headline UDE forecast accuracy. Full-state observation, 358 forecast points (t > 30 ms) of one uninterrupted rollout from t=0 — the model is never re-initialized at t_train_end, so this is extrapolation error, not one-step error. Scored against the CLEAN trajectory. Pair with UdeTrainVRmseVsNoisy for the noise floor.
% source=results/metrics_baseline_multiseed.csv filter={model=UDE, observed=full, window=forecast, metric=V_rmse} column=mean/std raw={mean=0.24523933188134714, std=0.028375840021892113} n=5
\newcommand{\valUdeForecastVRmseMeanStd}{0.245 \pm 0.028}

% note: THE headline UDE forecast accuracy. Full-state observation, 358 forecast points (t > 30 ms) of one uninterrupted rollout from t=0 — the model is never re-initialized at t_train_end, so this is extrapolation error, not one-step error. Scored against the CLEAN trajectory. Pair with UdeTrainVRmseVsNoisy for the noise floor.
% source=results/metrics_baseline_multiseed.csv filter={model=UDE, observed=full, window=forecast, metric=V_rmse} column=mean raw=0.24523933188134714 n=5
\newcommand{\valUdeForecastVRmseMean}{0.245}

% note: THE headline UDE forecast accuracy. Full-state observation, 358 forecast points (t > 30 ms) of one uninterrupted rollout from t=0 — the model is never re-initialized at t_train_end, so this is extrapolation error, not one-step error. Scored against the CLEAN trajectory. Pair with UdeTrainVRmseVsNoisy for the noise floor.
% source=results/metrics_baseline_multiseed.csv filter={model=UDE, observed=full, window=forecast, metric=V_rmse} column=std raw=0.028375840021892113 n=5
\newcommand{\valUdeForecastVRmseStd}{0.028}

% note: THE headline UDE forecast accuracy. Full-state observation, 358 forecast points (t > 30 ms) of one uninterrupted rollout from t=0 — the model is never re-initialized at t_train_end, so this is extrapolation error, not one-step error. Scored against the CLEAN trajectory. Pair with UdeTrainVRmseVsNoisy for the noise floor.
% source=results/metrics_baseline_multiseed.csv filter={model=UDE, observed=full, window=forecast, metric=V_rmse} column=n_seeds raw=5 n=5
\newcommand{\valUdeForecastVRmseN}{5}

% note: Voltage-only observation: the loss selects row 1 only (make_ude_loss, experiment.jl:129-144); the gates are hidden from the OBJECTIVE, not from the solver (u0 is the full true 6-state IC in every run). Say "voltage-only observation in the loss", never "only V is known". EQUIVALENCE CLAIM ONLY — see claim-guards.md: 0.264 vs 0.245 against stds of 0.118/0.028 at n=5, with no test run anywhere in the codebase. Write "matches full-state observation to within seed noise". Do NOT write "voltage-only is slightly worse". Note also there are NO common_eval rows for observed=voltage, so that comparison cannot be made at all.
% source=results/metrics_baseline_multiseed.csv filter={model=UDE, observed=voltage, window=forecast, metric=V_rmse} column=mean/std raw={mean=0.2635508680094284, std=0.11760231838630762} n=5
\newcommand{\valUdeVoltageForecastVRmse}{0.264 \pm 0.118}

% note: Voltage-only observation: the loss selects row 1 only (make_ude_loss, experiment.jl:129-144); the gates are hidden from the OBJECTIVE, not from the solver (u0 is the full true 6-state IC in every run). Say "voltage-only observation in the loss", never "only V is known". EQUIVALENCE CLAIM ONLY — see claim-guards.md: 0.264 vs 0.245 against stds of 0.118/0.028 at n=5, with no test run anywhere in the codebase. Write "matches full-state observation to within seed noise". Do NOT write "voltage-only is slightly worse". Note also there are NO common_eval rows for observed=voltage, so that comparison cannot be made at all.
% source=results/metrics_baseline_multiseed.csv filter={model=UDE, observed=voltage, window=forecast, metric=V_rmse} column=mean/std raw={mean=0.2635508680094284, std=0.11760231838630762} n=5
\newcommand{\valUdeVoltageForecastVRmseMeanStd}{0.264 \pm 0.118}

% note: Voltage-only observation: the loss selects row 1 only (make_ude_loss, experiment.jl:129-144); the gates are hidden from the OBJECTIVE, not from the solver (u0 is the full true 6-state IC in every run). Say "voltage-only observation in the loss", never "only V is known". EQUIVALENCE CLAIM ONLY — see claim-guards.md: 0.264 vs 0.245 against stds of 0.118/0.028 at n=5, with no test run anywhere in the codebase. Write "matches full-state observation to within seed noise". Do NOT write "voltage-only is slightly worse". Note also there are NO common_eval rows for observed=voltage, so that comparison cannot be made at all.
% source=results/metrics_baseline_multiseed.csv filter={model=UDE, observed=voltage, window=forecast, metric=V_rmse} column=mean raw=0.2635508680094284 n=5
\newcommand{\valUdeVoltageForecastVRmseMean}{0.264}

% note: Voltage-only observation: the loss selects row 1 only (make_ude_loss, experiment.jl:129-144); the gates are hidden from the OBJECTIVE, not from the solver (u0 is the full true 6-state IC in every run). Say "voltage-only observation in the loss", never "only V is known". EQUIVALENCE CLAIM ONLY — see claim-guards.md: 0.264 vs 0.245 against stds of 0.118/0.028 at n=5, with no test run anywhere in the codebase. Write "matches full-state observation to within seed noise". Do NOT write "voltage-only is slightly worse". Note also there are NO common_eval rows for observed=voltage, so that comparison cannot be made at all.
% source=results/metrics_baseline_multiseed.csv filter={model=UDE, observed=voltage, window=forecast, metric=V_rmse} column=std raw=0.11760231838630762 n=5
\newcommand{\valUdeVoltageForecastVRmseStd}{0.118}

% note: Voltage-only observation: the loss selects row 1 only (make_ude_loss, experiment.jl:129-144); the gates are hidden from the OBJECTIVE, not from the solver (u0 is the full true 6-state IC in every run). Say "voltage-only observation in the loss", never "only V is known". EQUIVALENCE CLAIM ONLY — see claim-guards.md: 0.264 vs 0.245 against stds of 0.118/0.028 at n=5, with no test run anywhere in the codebase. Write "matches full-state observation to within seed noise". Do NOT write "voltage-only is slightly worse". Note also there are NO common_eval rows for observed=voltage, so that comparison cannot be made at all.
% source=results/metrics_baseline_multiseed.csv filter={model=UDE, observed=voltage, window=forecast, metric=V_rmse} column=n_seeds raw=5 n=5
\newcommand{\valUdeVoltageForecastVRmseN}{5}

% note: Mean of the five per-gate RMSEs (m,h,n,p,s) — an arithmetic MEAN OF RMSEs, not a pooled RMSE (src/metrics.jl); each gate is weighted equally regardless of dynamic range. Describe it as such. round: 4 because the value is ~1e-3 and 3 dp would render it as 0.001 ± 0.000.
% source=results/metrics_baseline_multiseed.csv filter={model=UDE, observed=full, window=forecast, metric=gate_rmse_mean} column=mean/std raw={mean=0.000951766817970978, std=0.00017026484222479465} n=5
\newcommand{\valUdeForecastGateRmse}{0.0010 \pm 0.0002}

% note: Mean of the five per-gate RMSEs (m,h,n,p,s) — an arithmetic MEAN OF RMSEs, not a pooled RMSE (src/metrics.jl); each gate is weighted equally regardless of dynamic range. Describe it as such. round: 4 because the value is ~1e-3 and 3 dp would render it as 0.001 ± 0.000.
% source=results/metrics_baseline_multiseed.csv filter={model=UDE, observed=full, window=forecast, metric=gate_rmse_mean} column=mean/std raw={mean=0.000951766817970978, std=0.00017026484222479465} n=5
\newcommand{\valUdeForecastGateRmseMeanStd}{0.0010 \pm 0.0002}

% note: Mean of the five per-gate RMSEs (m,h,n,p,s) — an arithmetic MEAN OF RMSEs, not a pooled RMSE (src/metrics.jl); each gate is weighted equally regardless of dynamic range. Describe it as such. round: 4 because the value is ~1e-3 and 3 dp would render it as 0.001 ± 0.000.
% source=results/metrics_baseline_multiseed.csv filter={model=UDE, observed=full, window=forecast, metric=gate_rmse_mean} column=mean raw=0.000951766817970978 n=5
\newcommand{\valUdeForecastGateRmseMean}{0.0010}

% note: Mean of the five per-gate RMSEs (m,h,n,p,s) — an arithmetic MEAN OF RMSEs, not a pooled RMSE (src/metrics.jl); each gate is weighted equally regardless of dynamic range. Describe it as such. round: 4 because the value is ~1e-3 and 3 dp would render it as 0.001 ± 0.000.
% source=results/metrics_baseline_multiseed.csv filter={model=UDE, observed=full, window=forecast, metric=gate_rmse_mean} column=std raw=0.00017026484222479465 n=5
\newcommand{\valUdeForecastGateRmseStd}{0.0002}

% note: Mean of the five per-gate RMSEs (m,h,n,p,s) — an arithmetic MEAN OF RMSEs, not a pooled RMSE (src/metrics.jl); each gate is weighted equally regardless of dynamic range. Describe it as such. round: 4 because the value is ~1e-3 and 3 dp would render it as 0.001 ± 0.000.
% source=results/metrics_baseline_multiseed.csv filter={model=UDE, observed=full, window=forecast, metric=gate_rmse_mean} column=n_seeds raw=5 n=5
\newcommand{\valUdeForecastGateRmseN}{5}

% note: Hidden-calcium-current recovery error, µA/cm². UDE ONLY — ICa_* rows do not exist for NODE (no calcium closure), so there is no baseline comparison for this metric and none should be implied. CRITICAL: this is a FUNCTION-APPROXIMATION error, not a rollout error. ICa_nn is probed at the true CLEAN (V,s) (experiment.jl:363-365), decoupled from trajectory error, whereas V_rmse comes from a free 100 ms rollout. Different objects — do not narrate them as one.
% source=results/metrics_baseline_multiseed.csv filter={model=UDE, observed=full, window=forecast, metric=ICa_rmse} column=mean/std raw={mean=1.189072098980375, std=0.20829298482831607} n=5
\newcommand{\valUdeForecastICaRmse}{1.189 \pm 0.208}

% note: Hidden-calcium-current recovery error, µA/cm². UDE ONLY — ICa_* rows do not exist for NODE (no calcium closure), so there is no baseline comparison for this metric and none should be implied. CRITICAL: this is a FUNCTION-APPROXIMATION error, not a rollout error. ICa_nn is probed at the true CLEAN (V,s) (experiment.jl:363-365), decoupled from trajectory error, whereas V_rmse comes from a free 100 ms rollout. Different objects — do not narrate them as one.
% source=results/metrics_baseline_multiseed.csv filter={model=UDE, observed=full, window=forecast, metric=ICa_rmse} column=mean/std raw={mean=1.189072098980375, std=0.20829298482831607} n=5
\newcommand{\valUdeForecastICaRmseMeanStd}{1.189 \pm 0.208}

% note: Hidden-calcium-current recovery error, µA/cm². UDE ONLY — ICa_* rows do not exist for NODE (no calcium closure), so there is no baseline comparison for this metric and none should be implied. CRITICAL: this is a FUNCTION-APPROXIMATION error, not a rollout error. ICa_nn is probed at the true CLEAN (V,s) (experiment.jl:363-365), decoupled from trajectory error, whereas V_rmse comes from a free 100 ms rollout. Different objects — do not narrate them as one.
% source=results/metrics_baseline_multiseed.csv filter={model=UDE, observed=full, window=forecast, metric=ICa_rmse} column=mean raw=1.189072098980375 n=5
\newcommand{\valUdeForecastICaRmseMean}{1.189}

% note: Hidden-calcium-current recovery error, µA/cm². UDE ONLY — ICa_* rows do not exist for NODE (no calcium closure), so there is no baseline comparison for this metric and none should be implied. CRITICAL: this is a FUNCTION-APPROXIMATION error, not a rollout error. ICa_nn is probed at the true CLEAN (V,s) (experiment.jl:363-365), decoupled from trajectory error, whereas V_rmse comes from a free 100 ms rollout. Different objects — do not narrate them as one.
% source=results/metrics_baseline_multiseed.csv filter={model=UDE, observed=full, window=forecast, metric=ICa_rmse} column=std raw=0.20829298482831607 n=5
\newcommand{\valUdeForecastICaRmseStd}{0.208}

% note: Hidden-calcium-current recovery error, µA/cm². UDE ONLY — ICa_* rows do not exist for NODE (no calcium closure), so there is no baseline comparison for this metric and none should be implied. CRITICAL: this is a FUNCTION-APPROXIMATION error, not a rollout error. ICa_nn is probed at the true CLEAN (V,s) (experiment.jl:363-365), decoupled from trajectory error, whereas V_rmse comes from a free 100 ms rollout. Different objects — do not narrate them as one.
% source=results/metrics_baseline_multiseed.csv filter={model=UDE, observed=full, window=forecast, metric=ICa_rmse} column=n_seeds raw=5 n=5
\newcommand{\valUdeForecastICaRmseN}{5}

% note: Voltage-only twin of UdeForecastICaRmse. Equivalence claim only: Δ = 0.12 µA/cm² against a pooled std ≈ 0.35 at n=5. See the guard on UdeVoltageForecastVRmse.
% source=results/metrics_baseline_multiseed.csv filter={model=UDE, observed=voltage, window=forecast, metric=ICa_rmse} column=mean/std raw={mean=1.3092265677039077, std=0.44677450570074734} n=5
\newcommand{\valUdeVoltageForecastICaRmse}{1.309 \pm 0.447}

% note: Voltage-only twin of UdeForecastICaRmse. Equivalence claim only: Δ = 0.12 µA/cm² against a pooled std ≈ 0.35 at n=5. See the guard on UdeVoltageForecastVRmse.
% source=results/metrics_baseline_multiseed.csv filter={model=UDE, observed=voltage, window=forecast, metric=ICa_rmse} column=mean/std raw={mean=1.3092265677039077, std=0.44677450570074734} n=5
\newcommand{\valUdeVoltageForecastICaRmseMeanStd}{1.309 \pm 0.447}

% note: Voltage-only twin of UdeForecastICaRmse. Equivalence claim only: Δ = 0.12 µA/cm² against a pooled std ≈ 0.35 at n=5. See the guard on UdeVoltageForecastVRmse.
% source=results/metrics_baseline_multiseed.csv filter={model=UDE, observed=voltage, window=forecast, metric=ICa_rmse} column=mean raw=1.3092265677039077 n=5
\newcommand{\valUdeVoltageForecastICaRmseMean}{1.309}

% note: Voltage-only twin of UdeForecastICaRmse. Equivalence claim only: Δ = 0.12 µA/cm² against a pooled std ≈ 0.35 at n=5. See the guard on UdeVoltageForecastVRmse.
% source=results/metrics_baseline_multiseed.csv filter={model=UDE, observed=voltage, window=forecast, metric=ICa_rmse} column=std raw=0.44677450570074734 n=5
\newcommand{\valUdeVoltageForecastICaRmseStd}{0.447}

% note: Voltage-only twin of UdeForecastICaRmse. Equivalence claim only: Δ = 0.12 µA/cm² against a pooled std ≈ 0.35 at n=5. See the guard on UdeVoltageForecastVRmse.
% source=results/metrics_baseline_multiseed.csv filter={model=UDE, observed=voltage, window=forecast, metric=ICa_rmse} column=n_seeds raw=5 n=5
\newcommand{\valUdeVoltageForecastICaRmseN}{5}

% note: ICa_rmse / std(ICa_true) WITHIN THAT WINDOW (src/metrics.jl:151-153). The normalizer is recomputed per window, so the train and forecast values use different denominators and are not comparable to each other. Its value is as the unit-free scale that makes the gCa sweep (Block 3) interpretable: >1 means the closure's error exceeds the signal's own variation. This is also the seed-spread yardstick (±0.025) for judging which single-seed ablation gaps in Block 3 are real.
% source=results/metrics_baseline_multiseed.csv filter={model=UDE, observed=full, window=forecast, metric=ICa_rmse_norm} column=mean/std raw={mean=0.1403411215254637, std=0.02458393492014376} n=5
\newcommand{\valUdeForecastICaRmseNorm}{0.140 \pm 0.025}

% note: ICa_rmse / std(ICa_true) WITHIN THAT WINDOW (src/metrics.jl:151-153). The normalizer is recomputed per window, so the train and forecast values use different denominators and are not comparable to each other. Its value is as the unit-free scale that makes the gCa sweep (Block 3) interpretable: >1 means the closure's error exceeds the signal's own variation. This is also the seed-spread yardstick (±0.025) for judging which single-seed ablation gaps in Block 3 are real.
% source=results/metrics_baseline_multiseed.csv filter={model=UDE, observed=full, window=forecast, metric=ICa_rmse_norm} column=mean/std raw={mean=0.1403411215254637, std=0.02458393492014376} n=5
\newcommand{\valUdeForecastICaRmseNormMeanStd}{0.140 \pm 0.025}

% note: ICa_rmse / std(ICa_true) WITHIN THAT WINDOW (src/metrics.jl:151-153). The normalizer is recomputed per window, so the train and forecast values use different denominators and are not comparable to each other. Its value is as the unit-free scale that makes the gCa sweep (Block 3) interpretable: >1 means the closure's error exceeds the signal's own variation. This is also the seed-spread yardstick (±0.025) for judging which single-seed ablation gaps in Block 3 are real.
% source=results/metrics_baseline_multiseed.csv filter={model=UDE, observed=full, window=forecast, metric=ICa_rmse_norm} column=mean raw=0.1403411215254637 n=5
\newcommand{\valUdeForecastICaRmseNormMean}{0.140}

% note: ICa_rmse / std(ICa_true) WITHIN THAT WINDOW (src/metrics.jl:151-153). The normalizer is recomputed per window, so the train and forecast values use different denominators and are not comparable to each other. Its value is as the unit-free scale that makes the gCa sweep (Block 3) interpretable: >1 means the closure's error exceeds the signal's own variation. This is also the seed-spread yardstick (±0.025) for judging which single-seed ablation gaps in Block 3 are real.
% source=results/metrics_baseline_multiseed.csv filter={model=UDE, observed=full, window=forecast, metric=ICa_rmse_norm} column=std raw=0.02458393492014376 n=5
\newcommand{\valUdeForecastICaRmseNormStd}{0.025}

% note: ICa_rmse / std(ICa_true) WITHIN THAT WINDOW (src/metrics.jl:151-153). The normalizer is recomputed per window, so the train and forecast values use different denominators and are not comparable to each other. Its value is as the unit-free scale that makes the gCa sweep (Block 3) interpretable: >1 means the closure's error exceeds the signal's own variation. This is also the seed-spread yardstick (±0.025) for judging which single-seed ablation gaps in Block 3 are real.
% source=results/metrics_baseline_multiseed.csv filter={model=UDE, observed=full, window=forecast, metric=ICa_rmse_norm} column=n_seeds raw=5 n=5
\newcommand{\valUdeForecastICaRmseNormN}{5}

% note: Voltage-only twin of UdeForecastICaRmseNorm. Equivalence claim only.
% source=results/metrics_baseline_multiseed.csv filter={model=UDE, observed=voltage, window=forecast, metric=ICa_rmse_norm} column=mean/std raw={mean=0.15452244233133955, std=0.05273089432742834} n=5
\newcommand{\valUdeVoltageForecastICaRmseNorm}{0.155 \pm 0.053}

% note: Voltage-only twin of UdeForecastICaRmseNorm. Equivalence claim only.
% source=results/metrics_baseline_multiseed.csv filter={model=UDE, observed=voltage, window=forecast, metric=ICa_rmse_norm} column=mean/std raw={mean=0.15452244233133955, std=0.05273089432742834} n=5
\newcommand{\valUdeVoltageForecastICaRmseNormMeanStd}{0.155 \pm 0.053}

% note: Voltage-only twin of UdeForecastICaRmseNorm. Equivalence claim only.
% source=results/metrics_baseline_multiseed.csv filter={model=UDE, observed=voltage, window=forecast, metric=ICa_rmse_norm} column=mean raw=0.15452244233133955 n=5
\newcommand{\valUdeVoltageForecastICaRmseNormMean}{0.155}

% note: Voltage-only twin of UdeForecastICaRmseNorm. Equivalence claim only.
% source=results/metrics_baseline_multiseed.csv filter={model=UDE, observed=voltage, window=forecast, metric=ICa_rmse_norm} column=std raw=0.05273089432742834 n=5
\newcommand{\valUdeVoltageForecastICaRmseNormStd}{0.053}

% note: Voltage-only twin of UdeForecastICaRmseNorm. Equivalence claim only.
% source=results/metrics_baseline_multiseed.csv filter={model=UDE, observed=voltage, window=forecast, metric=ICa_rmse_norm} column=n_seeds raw=5 n=5
\newcommand{\valUdeVoltageForecastICaRmseNormN}{5}

% note: Train-window (t ≤ 30 ms, 154 points) V RMSE vs the CLEAN trajectory. Quote beside UdeForecastVRmse to show the train→forecast gap is small (0.148 → 0.245) — the UDE is not memorizing the window. Contrast NODE, which is 9.45 ± 3.21 on train and diverges on forecast.
% source=results/metrics_baseline_multiseed.csv filter={model=UDE, observed=full, window=train, metric=V_rmse} column=mean/std raw={mean=0.14750473251115578, std=0.04544861830300034} n=5
\newcommand{\valUdeTrainVRmse}{0.148 \pm 0.045}

% note: Train-window (t ≤ 30 ms, 154 points) V RMSE vs the CLEAN trajectory. Quote beside UdeForecastVRmse to show the train→forecast gap is small (0.148 → 0.245) — the UDE is not memorizing the window. Contrast NODE, which is 9.45 ± 3.21 on train and diverges on forecast.
% source=results/metrics_baseline_multiseed.csv filter={model=UDE, observed=full, window=train, metric=V_rmse} column=mean/std raw={mean=0.14750473251115578, std=0.04544861830300034} n=5
\newcommand{\valUdeTrainVRmseMeanStd}{0.148 \pm 0.045}

% note: Train-window (t ≤ 30 ms, 154 points) V RMSE vs the CLEAN trajectory. Quote beside UdeForecastVRmse to show the train→forecast gap is small (0.148 → 0.245) — the UDE is not memorizing the window. Contrast NODE, which is 9.45 ± 3.21 on train and diverges on forecast.
% source=results/metrics_baseline_multiseed.csv filter={model=UDE, observed=full, window=train, metric=V_rmse} column=mean raw=0.14750473251115578 n=5
\newcommand{\valUdeTrainVRmseMean}{0.148}

% note: Train-window (t ≤ 30 ms, 154 points) V RMSE vs the CLEAN trajectory. Quote beside UdeForecastVRmse to show the train→forecast gap is small (0.148 → 0.245) — the UDE is not memorizing the window. Contrast NODE, which is 9.45 ± 3.21 on train and diverges on forecast.
% source=results/metrics_baseline_multiseed.csv filter={model=UDE, observed=full, window=train, metric=V_rmse} column=std raw=0.04544861830300034 n=5
\newcommand{\valUdeTrainVRmseStd}{0.045}

% note: Train-window (t ≤ 30 ms, 154 points) V RMSE vs the CLEAN trajectory. Quote beside UdeForecastVRmse to show the train→forecast gap is small (0.148 → 0.245) — the UDE is not memorizing the window. Contrast NODE, which is 9.45 ± 3.21 on train and diverges on forecast.
% source=results/metrics_baseline_multiseed.csv filter={model=UDE, observed=full, window=train, metric=V_rmse} column=n_seeds raw=5 n=5
\newcommand{\valUdeTrainVRmseN}{5}

% note: THE NOISE FLOOR — the load-bearing companion to UdeTrainVRmse. Same prediction, scored against the NOISY training target instead of the clean truth. A model that recovered the clean signal exactly would score ≈ the noise std here, not 0. That UdeTrainVRmse (0.148, vs clean) is well BELOW UdeTrainVRmseVsNoisy (0.481, vs noisy) is the direct evidence the UDE fit the signal and not the noise. Without this macro the reader cannot tell whether 0.148 mV is good.
% source=results/metrics_baseline_multiseed.csv filter={model=UDE, observed=full, window=train, metric=V_rmse_vs_noisy} column=mean/std raw={mean=0.48134884129865074, std=0.012359910532320035} n=5
\newcommand{\valUdeTrainVRmseVsNoisy}{0.481 \pm 0.012}

% note: THE NOISE FLOOR — the load-bearing companion to UdeTrainVRmse. Same prediction, scored against the NOISY training target instead of the clean truth. A model that recovered the clean signal exactly would score ≈ the noise std here, not 0. That UdeTrainVRmse (0.148, vs clean) is well BELOW UdeTrainVRmseVsNoisy (0.481, vs noisy) is the direct evidence the UDE fit the signal and not the noise. Without this macro the reader cannot tell whether 0.148 mV is good.
% source=results/metrics_baseline_multiseed.csv filter={model=UDE, observed=full, window=train, metric=V_rmse_vs_noisy} column=mean/std raw={mean=0.48134884129865074, std=0.012359910532320035} n=5
\newcommand{\valUdeTrainVRmseVsNoisyMeanStd}{0.481 \pm 0.012}

% note: THE NOISE FLOOR — the load-bearing companion to UdeTrainVRmse. Same prediction, scored against the NOISY training target instead of the clean truth. A model that recovered the clean signal exactly would score ≈ the noise std here, not 0. That UdeTrainVRmse (0.148, vs clean) is well BELOW UdeTrainVRmseVsNoisy (0.481, vs noisy) is the direct evidence the UDE fit the signal and not the noise. Without this macro the reader cannot tell whether 0.148 mV is good.
% source=results/metrics_baseline_multiseed.csv filter={model=UDE, observed=full, window=train, metric=V_rmse_vs_noisy} column=mean raw=0.48134884129865074 n=5
\newcommand{\valUdeTrainVRmseVsNoisyMean}{0.481}

% note: THE NOISE FLOOR — the load-bearing companion to UdeTrainVRmse. Same prediction, scored against the NOISY training target instead of the clean truth. A model that recovered the clean signal exactly would score ≈ the noise std here, not 0. That UdeTrainVRmse (0.148, vs clean) is well BELOW UdeTrainVRmseVsNoisy (0.481, vs noisy) is the direct evidence the UDE fit the signal and not the noise. Without this macro the reader cannot tell whether 0.148 mV is good.
% source=results/metrics_baseline_multiseed.csv filter={model=UDE, observed=full, window=train, metric=V_rmse_vs_noisy} column=std raw=0.012359910532320035 n=5
\newcommand{\valUdeTrainVRmseVsNoisyStd}{0.012}

% note: THE NOISE FLOOR — the load-bearing companion to UdeTrainVRmse. Same prediction, scored against the NOISY training target instead of the clean truth. A model that recovered the clean signal exactly would score ≈ the noise std here, not 0. That UdeTrainVRmse (0.148, vs clean) is well BELOW UdeTrainVRmseVsNoisy (0.481, vs noisy) is the direct evidence the UDE fit the signal and not the noise. Without this macro the reader cannot tell whether 0.148 mV is good.
% source=results/metrics_baseline_multiseed.csv filter={model=UDE, observed=full, window=train, metric=V_rmse_vs_noisy} column=n_seeds raw=5 n=5
\newcommand{\valUdeTrainVRmseVsNoisyN}{5}

% note: *** THE HONEST NODE HEADLINE — use this, not NodeForecastVRmse (G2). *** Milliseconds until |V_pred − V_true| exceeds 10 mV for a sustained ≥1 ms. dt ≈ 0.196 ms, so 0.55 ms means the black-box Neural ODE leaves the tube within ONE TO EIGHT SAMPLES; per-seed 0.783/0.0/0.391/1.566/0.0, and the two 0.0 seeds breached at the very first forecast sample. Uncensored (far from the 70 ms cap), physically interpretable, and it states the failure without the spurious precision of a diverged RMSE. Not censored ⇒ mean_std is legitimate here, unlike UdeForecastRollout.
% source=results/metrics_baseline_multiseed.csv filter={model=NODE, observed=full, window=forecast, metric=rollout_horizon_ms} column=mean/std raw={mean=0.5479452054794521, std=0.6549197859366547} n=5
\newcommand{\valNodeForecastRollout}{0.55 \pm 0.65}

% note: *** THE HONEST NODE HEADLINE — use this, not NodeForecastVRmse (G2). *** Milliseconds until |V_pred − V_true| exceeds 10 mV for a sustained ≥1 ms. dt ≈ 0.196 ms, so 0.55 ms means the black-box Neural ODE leaves the tube within ONE TO EIGHT SAMPLES; per-seed 0.783/0.0/0.391/1.566/0.0, and the two 0.0 seeds breached at the very first forecast sample. Uncensored (far from the 70 ms cap), physically interpretable, and it states the failure without the spurious precision of a diverged RMSE. Not censored ⇒ mean_std is legitimate here, unlike UdeForecastRollout.
% source=results/metrics_baseline_multiseed.csv filter={model=NODE, observed=full, window=forecast, metric=rollout_horizon_ms} column=mean/std raw={mean=0.5479452054794521, std=0.6549197859366547} n=5
\newcommand{\valNodeForecastRolloutMeanStd}{0.55 \pm 0.65}

% note: *** THE HONEST NODE HEADLINE — use this, not NodeForecastVRmse (G2). *** Milliseconds until |V_pred − V_true| exceeds 10 mV for a sustained ≥1 ms. dt ≈ 0.196 ms, so 0.55 ms means the black-box Neural ODE leaves the tube within ONE TO EIGHT SAMPLES; per-seed 0.783/0.0/0.391/1.566/0.0, and the two 0.0 seeds breached at the very first forecast sample. Uncensored (far from the 70 ms cap), physically interpretable, and it states the failure without the spurious precision of a diverged RMSE. Not censored ⇒ mean_std is legitimate here, unlike UdeForecastRollout.
% source=results/metrics_baseline_multiseed.csv filter={model=NODE, observed=full, window=forecast, metric=rollout_horizon_ms} column=mean raw=0.5479452054794521 n=5
\newcommand{\valNodeForecastRolloutMean}{0.55}

% note: *** THE HONEST NODE HEADLINE — use this, not NodeForecastVRmse (G2). *** Milliseconds until |V_pred − V_true| exceeds 10 mV for a sustained ≥1 ms. dt ≈ 0.196 ms, so 0.55 ms means the black-box Neural ODE leaves the tube within ONE TO EIGHT SAMPLES; per-seed 0.783/0.0/0.391/1.566/0.0, and the two 0.0 seeds breached at the very first forecast sample. Uncensored (far from the 70 ms cap), physically interpretable, and it states the failure without the spurious precision of a diverged RMSE. Not censored ⇒ mean_std is legitimate here, unlike UdeForecastRollout.
% source=results/metrics_baseline_multiseed.csv filter={model=NODE, observed=full, window=forecast, metric=rollout_horizon_ms} column=std raw=0.6549197859366547 n=5
\newcommand{\valNodeForecastRolloutStd}{0.65}

% note: *** THE HONEST NODE HEADLINE — use this, not NodeForecastVRmse (G2). *** Milliseconds until |V_pred − V_true| exceeds 10 mV for a sustained ≥1 ms. dt ≈ 0.196 ms, so 0.55 ms means the black-box Neural ODE leaves the tube within ONE TO EIGHT SAMPLES; per-seed 0.783/0.0/0.391/1.566/0.0, and the two 0.0 seeds breached at the very first forecast sample. Uncensored (far from the 70 ms cap), physically interpretable, and it states the failure without the spurious precision of a diverged RMSE. Not censored ⇒ mean_std is legitimate here, unlike UdeForecastRollout.
% source=results/metrics_baseline_multiseed.csv filter={model=NODE, observed=full, window=forecast, metric=rollout_horizon_ms} column=n_seeds raw=5 n=5
\newcommand{\valNodeForecastRolloutN}{5}

% note: *** GUARDED — claim-guards.md G1. emit: ge IS MANDATORY HERE. *** 70.0 is cap_ms (src/metrics.jl:81-82), returned when the 10 mV tube is NEVER breached — a right-censored value, not a measured horizon. It is also exactly 100 − 30, i.e. the entire remaining trajectory: the metric is saturated for all 5 seeds and cannot discriminate among them. A macro emitting "70.0 ± 0.0 ms" asserts zero variance on a number that is really "≥ 70 and unknown" — the single most misleading rendering available in this results tree. Renders "\ge 70"; the caption must name the cap. Same censoring applies to UDE/voltage (also 70.0 ± 0.0); add a parallel entry only if the paper quotes it separately.
% source=results/metrics_baseline_multiseed.csv filter={model=UDE, observed=full, window=forecast, metric=rollout_horizon_ms} column=mean raw=70.0 n=1
\newcommand{\valUdeForecastRollout}{\ge 70}

% note: gCa = 0.4 (the physiological value). >1 means THE CLOSURE'S ERROR EXCEEDS THE TRUE CURRENT'S OWN STANDARD DEVIATION — the network has learned nothing about I_Ca at this conductance. This is the strongest single number in the identifiability argument and the reason the paper's main setting is the non-physiological gCa=2.0. Single seed (G13).
% source=results/ablation_gca.csv filter={gCa=0.4} column=forecast_ICa_rmse_norm raw=1.1791588919620495 n=1
\newcommand{\valGcaAblICaNormPointFour}{1.179}

% note: gCa = 1.0. Single seed (G13). The 0.4→1.0 gap clears the ±0.025 seed spread.
% source=results/ablation_gca.csv filter={gCa=1.0} column=forecast_ICa_rmse_norm raw=0.3325491936637111 n=1
\newcommand{\valGcaAblICaNormOne}{0.333}

% note: gCa = 2.0, the paper's main setting. Single seed (G13). Sanity check: this is seed 1111 of the same configuration whose 5-seed mean is UdeForecastICaRmseNorm = 0.140 ± 0.025 — consistent, and a useful cross-check that the ablation row and the multiseed row describe the same run. Do not present 0.112 and 0.140 as two different results.
% source=results/ablation_gca.csv filter={gCa=2.0} column=forecast_ICa_rmse_norm raw=0.11229216448147902 n=1
\newcommand{\valGcaAblICaNormTwo}{0.112}

% note: gCa = 4.0. Single seed (G13). The 2.0→4.0 improvement (0.112 → 0.068) is within seed noise — quote it as part of the monotone trend, never as a standalone difference.
% source=results/ablation_gca.csv filter={gCa=4.0} column=forecast_ICa_rmse_norm raw=0.06772530981692793 n=1
\newcommand{\valGcaAblICaNormFour}{0.068}

% note: t_train_end = 15 ms. The one window result that clears seed noise: ~7x worse than 20 ms. Claim: below ~20 ms of data the UDE fails to constrain the closure. Guards G11 (commoneval), G12 (artifact caveat), G13 (n=1).
% source=results/ablation_window.csv filter={t_train_end=15.0} column=common_eval_V_rmse raw=1.5586033539127329 n=1
\newcommand{\valWindowAblCommonEvalVRmseFifteen}{1.559}

% note: t_train_end = 20 ms. Guards G11, G12, G13. Gap to 30 ms is within seed noise.
% source=results/ablation_window.csv filter={t_train_end=20.0} column=common_eval_V_rmse raw=0.22523666920122537 n=1
\newcommand{\valWindowAblCommonEvalVRmseTwenty}{0.225}

% note: t_train_end = 30 ms — the baseline window, seed 1111. Guards G11, G12, G13. Not the same quantity as UdeForecastVRmse (0.245 ± 0.028): that is the forecast window (t > 30) over 5 seeds, this is the common-eval window (t > 50) at one seed. Do not conflate them.
% source=results/ablation_window.csv filter={t_train_end=30.0} column=common_eval_V_rmse raw=0.23608373353853263 n=1
\newcommand{\valWindowAblCommonEvalVRmseThirty}{0.236}

% note: t_train_end = 40 ms. *** Worse than 30 ms — G12 artifact, MUST be flagged wherever quoted. *** More training data with a fixed iteration budget = under-optimisation, not a harder problem.
% source=results/ablation_window.csv filter={t_train_end=40.0} column=common_eval_V_rmse raw=0.3671703369134956 n=1
\newcommand{\valWindowAblCommonEvalVRmseForty}{0.367}

% note: t_train_end = 50 ms. G12 artifact (worst of the four non-collapsed windows). Additional caveat: at this point common_eval_* == forecast_* exactly, since make_split(50.0) and common_eval_indices() select the same index set — correct behaviour, but not an independent check.
% source=results/ablation_window.csv filter={t_train_end=50.0} column=common_eval_V_rmse raw=0.4177206252114012 n=1
\newcommand{\valWindowAblCommonEvalVRmseFifty}{0.418}

% note: The per-seed span of a_hat, and the number that makes SymbolicAHat's "1.60 ± 0.90" legible: the estimates span 4.6x and BRACKET the true 2.0 rather than concentrating near it. Text mode — renders an en-dash range, not math. A SECOND entry over the same column as SymbolicAHat is the correct construction: one entry has one `emit:`, and mean_std has already spent SymbolicAHat's. The suffix-collision worry is unfounded — only `emit: mean_std` appends Mean/Std/N/MeanStd; `emit: range` emits the bare name, so this entry generates exactly \valSymbolicAHatRange.
% source=results/symbolic/symbolic_recovery_metrics.csv filter={gCa=2.0} column=a_hat raw={min=0.6730558988367701, max=3.094842062703728} n=5
\newcommand{\valSymbolicAHatRange}{0.673--3.095}

% note: *** THE NEGATIVE RESULT. 1.60 ± 0.90 against a true 2.0 IS NOT A RECOVERY. *** The std is 57% of the mean and the per-seed values span 4.6x. Quote it ONLY as "order-correct", never as "recovers gCa". The exoneration control that makes this a claim about the data rather than the code: the same estimator fitted to the TRUE current on the same hull-masked domain returns a ≈ 2.0, b ≈ −240 (objective3_symbolic.jl:208-209 `[sanity]`) — so the fit machinery is sound and the spread is the NETWORK's non-identifiability. Cite that control whenever this macro appears. DO NOT also quote the "ensemble" a = 1.596 as a corroborating second estimate — it is algebraically the same number (see the DO-NOT-MACRO block).
% source=results/symbolic/symbolic_recovery_metrics.csv filter={gCa=2.0} column=a_hat (aggregated over seed) raw={mean=1.595549902182419, std=0.9045688330069698} n=5
\newcommand{\valSymbolicAHat}{1.60 \pm 0.90}

% note: *** THE NEGATIVE RESULT. 1.60 ± 0.90 against a true 2.0 IS NOT A RECOVERY. *** The std is 57% of the mean and the per-seed values span 4.6x. Quote it ONLY as "order-correct", never as "recovers gCa". The exoneration control that makes this a claim about the data rather than the code: the same estimator fitted to the TRUE current on the same hull-masked domain returns a ≈ 2.0, b ≈ −240 (objective3_symbolic.jl:208-209 `[sanity]`) — so the fit machinery is sound and the spread is the NETWORK's non-identifiability. Cite that control whenever this macro appears. DO NOT also quote the "ensemble" a = 1.596 as a corroborating second estimate — it is algebraically the same number (see the DO-NOT-MACRO block).
% source=results/symbolic/symbolic_recovery_metrics.csv filter={gCa=2.0} column=a_hat (aggregated over seed) raw={mean=1.595549902182419, std=0.9045688330069698} n=5
\newcommand{\valSymbolicAHatMeanStd}{1.60 \pm 0.90}

% note: *** THE NEGATIVE RESULT. 1.60 ± 0.90 against a true 2.0 IS NOT A RECOVERY. *** The std is 57% of the mean and the per-seed values span 4.6x. Quote it ONLY as "order-correct", never as "recovers gCa". The exoneration control that makes this a claim about the data rather than the code: the same estimator fitted to the TRUE current on the same hull-masked domain returns a ≈ 2.0, b ≈ −240 (objective3_symbolic.jl:208-209 `[sanity]`) — so the fit machinery is sound and the spread is the NETWORK's non-identifiability. Cite that control whenever this macro appears. DO NOT also quote the "ensemble" a = 1.596 as a corroborating second estimate — it is algebraically the same number (see the DO-NOT-MACRO block).
% source=results/symbolic/symbolic_recovery_metrics.csv filter={gCa=2.0} column=a_hat raw=1.595549902182419 n=5
\newcommand{\valSymbolicAHatMean}{1.60}

% note: *** THE NEGATIVE RESULT. 1.60 ± 0.90 against a true 2.0 IS NOT A RECOVERY. *** The std is 57% of the mean and the per-seed values span 4.6x. Quote it ONLY as "order-correct", never as "recovers gCa". The exoneration control that makes this a claim about the data rather than the code: the same estimator fitted to the TRUE current on the same hull-masked domain returns a ≈ 2.0, b ≈ −240 (objective3_symbolic.jl:208-209 `[sanity]`) — so the fit machinery is sound and the spread is the NETWORK's non-identifiability. Cite that control whenever this macro appears. DO NOT also quote the "ensemble" a = 1.596 as a corroborating second estimate — it is algebraically the same number (see the DO-NOT-MACRO block).
% source=results/symbolic/symbolic_recovery_metrics.csv filter={gCa=2.0} column=a_hat raw=0.9045688330069698 n=5
\newcommand{\valSymbolicAHatStd}{0.90}

% note: *** THE NEGATIVE RESULT. 1.60 ± 0.90 against a true 2.0 IS NOT A RECOVERY. *** The std is 57% of the mean and the per-seed values span 4.6x. Quote it ONLY as "order-correct", never as "recovers gCa". The exoneration control that makes this a claim about the data rather than the code: the same estimator fitted to the TRUE current on the same hull-masked domain returns a ≈ 2.0, b ≈ −240 (objective3_symbolic.jl:208-209 `[sanity]`) — so the fit machinery is sound and the spread is the NETWORK's non-identifiability. Cite that control whenever this macro appears. DO NOT also quote the "ensemble" a = 1.596 as a corroborating second estimate — it is algebraically the same number (see the DO-NOT-MACRO block).
% source=results/symbolic/symbolic_recovery_metrics.csv filter={gCa=2.0} column=seed raw=5 n=5
\newcommand{\valSymbolicAHatN}{5}

% note: b ↔ −gCa·ECa, true −240.0. Closer in relative terms than a_hat (8% off) but with a 23% std. Same "order-correct, not recovered" framing. round: 1 because the value is O(10²) — 3 dp would be absurd precision on ±59.5.
% source=results/symbolic/symbolic_recovery_metrics.csv filter={gCa=2.0} column=b_hat (aggregated over seed) raw={mean=-259.33239319241306, std=59.501608270568674} n=5
\newcommand{\valSymbolicBHat}{-259.3 \pm 59.5}

% note: b ↔ −gCa·ECa, true −240.0. Closer in relative terms than a_hat (8% off) but with a 23% std. Same "order-correct, not recovered" framing. round: 1 because the value is O(10²) — 3 dp would be absurd precision on ±59.5.
% source=results/symbolic/symbolic_recovery_metrics.csv filter={gCa=2.0} column=b_hat (aggregated over seed) raw={mean=-259.33239319241306, std=59.501608270568674} n=5
\newcommand{\valSymbolicBHatMeanStd}{-259.3 \pm 59.5}

% note: b ↔ −gCa·ECa, true −240.0. Closer in relative terms than a_hat (8% off) but with a 23% std. Same "order-correct, not recovered" framing. round: 1 because the value is O(10²) — 3 dp would be absurd precision on ±59.5.
% source=results/symbolic/symbolic_recovery_metrics.csv filter={gCa=2.0} column=b_hat raw=-259.33239319241306 n=5
\newcommand{\valSymbolicBHatMean}{-259.3}

% note: b ↔ −gCa·ECa, true −240.0. Closer in relative terms than a_hat (8% off) but with a 23% std. Same "order-correct, not recovered" framing. round: 1 because the value is O(10²) — 3 dp would be absurd precision on ±59.5.
% source=results/symbolic/symbolic_recovery_metrics.csv filter={gCa=2.0} column=b_hat raw=59.501608270568674 n=5
\newcommand{\valSymbolicBHatStd}{59.5}

% note: b ↔ −gCa·ECa, true −240.0. Closer in relative terms than a_hat (8% off) but with a 23% std. Same "order-correct, not recovered" framing. round: 1 because the value is O(10²) — 3 dp would be absurd precision on ±59.5.
% source=results/symbolic/symbolic_recovery_metrics.csv filter={gCa=2.0} column=seed raw=5 n=5
\newcommand{\valSymbolicBHatN}{5}

% note: *** THE R² THE PAPER MUST QUOTE — not R2_ens (see DO-NOT-MACRO block). *** R² of the conductance-form fit against the TRUE current, per seed, on the 512-point trajectory. This is the positive half of the paper's headline: the closure recovers the FUNCTIONAL FORM s²(aV+b) at R² = 0.984 ± 0.014 while its COEFFICIENTS (SymbolicAHat) do not identify. Form yes, coefficients no — state both in the same sentence or the result is misread. Note the domain: per-seed R² is on the trajectory; R2_ens is on the pooled hull grid. Not a like-for-like pair.
% source=results/symbolic/symbolic_recovery_metrics.csv filter={gCa=2.0} column=r2_cond_true (aggregated over seed) raw={mean=0.9835864514706014, std=0.014493538574901191} n=5
\newcommand{\valSymbolicRTwoCondTrue}{0.984 \pm 0.014}

% note: *** THE R² THE PAPER MUST QUOTE — not R2_ens (see DO-NOT-MACRO block). *** R² of the conductance-form fit against the TRUE current, per seed, on the 512-point trajectory. This is the positive half of the paper's headline: the closure recovers the FUNCTIONAL FORM s²(aV+b) at R² = 0.984 ± 0.014 while its COEFFICIENTS (SymbolicAHat) do not identify. Form yes, coefficients no — state both in the same sentence or the result is misread. Note the domain: per-seed R² is on the trajectory; R2_ens is on the pooled hull grid. Not a like-for-like pair.
% source=results/symbolic/symbolic_recovery_metrics.csv filter={gCa=2.0} column=r2_cond_true (aggregated over seed) raw={mean=0.9835864514706014, std=0.014493538574901191} n=5
\newcommand{\valSymbolicRTwoCondTrueMeanStd}{0.984 \pm 0.014}

% note: *** THE R² THE PAPER MUST QUOTE — not R2_ens (see DO-NOT-MACRO block). *** R² of the conductance-form fit against the TRUE current, per seed, on the 512-point trajectory. This is the positive half of the paper's headline: the closure recovers the FUNCTIONAL FORM s²(aV+b) at R² = 0.984 ± 0.014 while its COEFFICIENTS (SymbolicAHat) do not identify. Form yes, coefficients no — state both in the same sentence or the result is misread. Note the domain: per-seed R² is on the trajectory; R2_ens is on the pooled hull grid. Not a like-for-like pair.
% source=results/symbolic/symbolic_recovery_metrics.csv filter={gCa=2.0} column=r2_cond_true raw=0.9835864514706014 n=5
\newcommand{\valSymbolicRTwoCondTrueMean}{0.984}

% note: *** THE R² THE PAPER MUST QUOTE — not R2_ens (see DO-NOT-MACRO block). *** R² of the conductance-form fit against the TRUE current, per seed, on the 512-point trajectory. This is the positive half of the paper's headline: the closure recovers the FUNCTIONAL FORM s²(aV+b) at R² = 0.984 ± 0.014 while its COEFFICIENTS (SymbolicAHat) do not identify. Form yes, coefficients no — state both in the same sentence or the result is misread. Note the domain: per-seed R² is on the trajectory; R2_ens is on the pooled hull grid. Not a like-for-like pair.
% source=results/symbolic/symbolic_recovery_metrics.csv filter={gCa=2.0} column=r2_cond_true raw=0.014493538574901191 n=5
\newcommand{\valSymbolicRTwoCondTrueStd}{0.014}

% note: *** THE R² THE PAPER MUST QUOTE — not R2_ens (see DO-NOT-MACRO block). *** R² of the conductance-form fit against the TRUE current, per seed, on the 512-point trajectory. This is the positive half of the paper's headline: the closure recovers the FUNCTIONAL FORM s²(aV+b) at R² = 0.984 ± 0.014 while its COEFFICIENTS (SymbolicAHat) do not identify. Form yes, coefficients no — state both in the same sentence or the result is misread. Note the domain: per-seed R² is on the trajectory; R2_ens is on the pooled hull grid. Not a like-for-like pair.
% source=results/symbolic/symbolic_recovery_metrics.csv filter={gCa=2.0} column=seed raw=5 n=5
\newcommand{\valSymbolicRTwoCondTrueN}{5}

% note: Unconstrained STLSQ/SINDy on the 6-term library {1, V, s, V·s, s², V·s²} vs the true current. The paper's cleanest "predictive accuracy ≠ identifiability" illustration: R² = 0.956 while the fit NEVER SPARSIFIES (n_sindy_terms = 6/6 at every seed) and V·s² carries coefficient −1.469 at seed 1111 — THE WRONG SIGN against a true +2.0. High R², wrong physics. Contrast with SymbolicRTwoCondTrue, which is structure-informed (it hard-codes s² and the linear driving force) and so cannot fail structurally.
% source=results/symbolic/symbolic_recovery_metrics.csv filter={gCa=2.0} column=r2_sindy_true (aggregated over seed) raw={mean=0.9557821118990294, std=0.05915713448295907} n=5
\newcommand{\valSymbolicRTwoSindyTrue}{0.956 \pm 0.059}

% note: Unconstrained STLSQ/SINDy on the 6-term library {1, V, s, V·s, s², V·s²} vs the true current. The paper's cleanest "predictive accuracy ≠ identifiability" illustration: R² = 0.956 while the fit NEVER SPARSIFIES (n_sindy_terms = 6/6 at every seed) and V·s² carries coefficient −1.469 at seed 1111 — THE WRONG SIGN against a true +2.0. High R², wrong physics. Contrast with SymbolicRTwoCondTrue, which is structure-informed (it hard-codes s² and the linear driving force) and so cannot fail structurally.
% source=results/symbolic/symbolic_recovery_metrics.csv filter={gCa=2.0} column=r2_sindy_true (aggregated over seed) raw={mean=0.9557821118990294, std=0.05915713448295907} n=5
\newcommand{\valSymbolicRTwoSindyTrueMeanStd}{0.956 \pm 0.059}

% note: Unconstrained STLSQ/SINDy on the 6-term library {1, V, s, V·s, s², V·s²} vs the true current. The paper's cleanest "predictive accuracy ≠ identifiability" illustration: R² = 0.956 while the fit NEVER SPARSIFIES (n_sindy_terms = 6/6 at every seed) and V·s² carries coefficient −1.469 at seed 1111 — THE WRONG SIGN against a true +2.0. High R², wrong physics. Contrast with SymbolicRTwoCondTrue, which is structure-informed (it hard-codes s² and the linear driving force) and so cannot fail structurally.
% source=results/symbolic/symbolic_recovery_metrics.csv filter={gCa=2.0} column=r2_sindy_true raw=0.9557821118990294 n=5
\newcommand{\valSymbolicRTwoSindyTrueMean}{0.956}

% note: Unconstrained STLSQ/SINDy on the 6-term library {1, V, s, V·s, s², V·s²} vs the true current. The paper's cleanest "predictive accuracy ≠ identifiability" illustration: R² = 0.956 while the fit NEVER SPARSIFIES (n_sindy_terms = 6/6 at every seed) and V·s² carries coefficient −1.469 at seed 1111 — THE WRONG SIGN against a true +2.0. High R², wrong physics. Contrast with SymbolicRTwoCondTrue, which is structure-informed (it hard-codes s² and the linear driving force) and so cannot fail structurally.
% source=results/symbolic/symbolic_recovery_metrics.csv filter={gCa=2.0} column=r2_sindy_true raw=0.05915713448295907 n=5
\newcommand{\valSymbolicRTwoSindyTrueStd}{0.059}

% note: Unconstrained STLSQ/SINDy on the 6-term library {1, V, s, V·s, s², V·s²} vs the true current. The paper's cleanest "predictive accuracy ≠ identifiability" illustration: R² = 0.956 while the fit NEVER SPARSIFIES (n_sindy_terms = 6/6 at every seed) and V·s² carries coefficient −1.469 at seed 1111 — THE WRONG SIGN against a true +2.0. High R², wrong physics. Contrast with SymbolicRTwoCondTrue, which is structure-informed (it hard-codes s² and the linear driving force) and so cannot fail structurally.
% source=results/symbolic/symbolic_recovery_metrics.csv filter={gCa=2.0} column=seed raw=5 n=5
\newcommand{\valSymbolicRTwoSindyTrueN}{5}

% note: *** WRONG SIGN: −0.372 against a true +0.40. *** The negative control, and a far harder failure statement than its R² alone conveys — the recovered "conductance" is negative, i.e. the fit says calcium is an outward current. Companion: Erev = −239.87 mV against a true +120.0 (in the CSV's `Erev` column; add a macro if the paper quotes it). Single row, n=1, no error bar — never render this as ±anything. Pairs with GcaAblICaNormPointFour: at the physiological conductance both the closure and its distillation fail, and they fail consistently.
% source=results/symbolic/symbolic_recovery_metrics.csv filter={gCa=0.4} column=a_hat raw=-0.3719798574860592 n=1
\newcommand{\valSymbolicAHatNegControl}{-0.372}

% note: Negative control R² for the conductance form: 0.781 vs 0.984 at gCa=2.0 (SymbolicRTwoCondTrue). Drives fig12's two-panel contrast. Caution: 0.781 still sounds respectable — always quote it beside SymbolicAHatNegControl's sign error, which is what actually demonstrates the failure. n=1.
% source=results/symbolic/symbolic_recovery_metrics.csv filter={gCa=0.4} column=r2_cond_true raw=0.7806472493640154 n=1
\newcommand{\valSymbolicRTwoCondTrueNegControl}{0.781}

% note: *** NEGATIVE R²: −0.359 — the unconstrained SINDy fit is worse than predicting the mean. *** The unambiguous failure number at gCa=0.4, and the sharpest contrast in the paper (0.956 → −0.359 vs 0.984 → 0.781 for the structure-informed fit). It also shows what the conductance form's hard-coded s² is buying: structure props up a fit that carries no information. n=1.
% source=results/symbolic/symbolic_recovery_metrics.csv filter={gCa=0.4} column=r2_sindy_true raw=-0.35870485913710626 n=1
\newcommand{\valSymbolicRTwoSindyTrueNegControl}{-0.359}

% note: AFTER, vs SymbolicAHat BEFORE: the mean moved AWAY from the true 2.0 (1.60 → 1.49) and the std INFLATED 64% (0.90 → 1.48). One seed (5555) went sign-negative (a_hat = −0.877). 4x the iterations bought a worse and noisier estimate. Quote the pair, always in that order, always with RetrainFinalLoss to preempt "you just under-trained it".
% source=results/retrain_gca2_20k/symbolic/symbolic_recovery_metrics.csv filter={gCa=2.0} column=a_hat (aggregated over seed) raw={mean=1.4947096308576968, std=1.4842865354496968} n=5
\newcommand{\valRetrainAHat}{1.49 \pm 1.48}

% note: AFTER, vs SymbolicAHat BEFORE: the mean moved AWAY from the true 2.0 (1.60 → 1.49) and the std INFLATED 64% (0.90 → 1.48). One seed (5555) went sign-negative (a_hat = −0.877). 4x the iterations bought a worse and noisier estimate. Quote the pair, always in that order, always with RetrainFinalLoss to preempt "you just under-trained it".
% source=results/retrain_gca2_20k/symbolic/symbolic_recovery_metrics.csv filter={gCa=2.0} column=a_hat (aggregated over seed) raw={mean=1.4947096308576968, std=1.4842865354496968} n=5
\newcommand{\valRetrainAHatMeanStd}{1.49 \pm 1.48}

% note: AFTER, vs SymbolicAHat BEFORE: the mean moved AWAY from the true 2.0 (1.60 → 1.49) and the std INFLATED 64% (0.90 → 1.48). One seed (5555) went sign-negative (a_hat = −0.877). 4x the iterations bought a worse and noisier estimate. Quote the pair, always in that order, always with RetrainFinalLoss to preempt "you just under-trained it".
% source=results/retrain_gca2_20k/symbolic/symbolic_recovery_metrics.csv filter={gCa=2.0} column=a_hat raw=1.4947096308576968 n=5
\newcommand{\valRetrainAHatMean}{1.49}

% note: AFTER, vs SymbolicAHat BEFORE: the mean moved AWAY from the true 2.0 (1.60 → 1.49) and the std INFLATED 64% (0.90 → 1.48). One seed (5555) went sign-negative (a_hat = −0.877). 4x the iterations bought a worse and noisier estimate. Quote the pair, always in that order, always with RetrainFinalLoss to preempt "you just under-trained it".
% source=results/retrain_gca2_20k/symbolic/symbolic_recovery_metrics.csv filter={gCa=2.0} column=a_hat raw=1.4842865354496968 n=5
\newcommand{\valRetrainAHatStd}{1.48}

% note: AFTER, vs SymbolicAHat BEFORE: the mean moved AWAY from the true 2.0 (1.60 → 1.49) and the std INFLATED 64% (0.90 → 1.48). One seed (5555) went sign-negative (a_hat = −0.877). 4x the iterations bought a worse and noisier estimate. Quote the pair, always in that order, always with RetrainFinalLoss to preempt "you just under-trained it".
% source=results/retrain_gca2_20k/symbolic/symbolic_recovery_metrics.csv filter={gCa=2.0} column=seed raw=5 n=5
\newcommand{\valRetrainAHatN}{5}

% note: AFTER, vs SymbolicRTwoCondTrue BEFORE: 0.984 ± 0.014 → 0.845 ± 0.160. The functional-form recovery — the paper's POSITIVE result — degrades under 4x optimisation, and its std grows 11x. This same before/after pair is the evidence that R2_ens is worthless (it moved 0.997333 → 0.997341 across this collapse); see the DO-NOT-MACRO block.
% source=results/retrain_gca2_20k/symbolic/symbolic_recovery_metrics.csv filter={gCa=2.0} column=r2_cond_true (aggregated over seed) raw={mean=0.8454749165905042, std=0.1602837711598065} n=5
\newcommand{\valRetrainRTwoCondTrue}{0.845 \pm 0.160}

% note: AFTER, vs SymbolicRTwoCondTrue BEFORE: 0.984 ± 0.014 → 0.845 ± 0.160. The functional-form recovery — the paper's POSITIVE result — degrades under 4x optimisation, and its std grows 11x. This same before/after pair is the evidence that R2_ens is worthless (it moved 0.997333 → 0.997341 across this collapse); see the DO-NOT-MACRO block.
% source=results/retrain_gca2_20k/symbolic/symbolic_recovery_metrics.csv filter={gCa=2.0} column=r2_cond_true (aggregated over seed) raw={mean=0.8454749165905042, std=0.1602837711598065} n=5
\newcommand{\valRetrainRTwoCondTrueMeanStd}{0.845 \pm 0.160}

% note: AFTER, vs SymbolicRTwoCondTrue BEFORE: 0.984 ± 0.014 → 0.845 ± 0.160. The functional-form recovery — the paper's POSITIVE result — degrades under 4x optimisation, and its std grows 11x. This same before/after pair is the evidence that R2_ens is worthless (it moved 0.997333 → 0.997341 across this collapse); see the DO-NOT-MACRO block.
% source=results/retrain_gca2_20k/symbolic/symbolic_recovery_metrics.csv filter={gCa=2.0} column=r2_cond_true raw=0.8454749165905042 n=5
\newcommand{\valRetrainRTwoCondTrueMean}{0.845}

% note: AFTER, vs SymbolicRTwoCondTrue BEFORE: 0.984 ± 0.014 → 0.845 ± 0.160. The functional-form recovery — the paper's POSITIVE result — degrades under 4x optimisation, and its std grows 11x. This same before/after pair is the evidence that R2_ens is worthless (it moved 0.997333 → 0.997341 across this collapse); see the DO-NOT-MACRO block.
% source=results/retrain_gca2_20k/symbolic/symbolic_recovery_metrics.csv filter={gCa=2.0} column=r2_cond_true raw=0.1602837711598065 n=5
\newcommand{\valRetrainRTwoCondTrueStd}{0.160}

% note: AFTER, vs SymbolicRTwoCondTrue BEFORE: 0.984 ± 0.014 → 0.845 ± 0.160. The functional-form recovery — the paper's POSITIVE result — degrades under 4x optimisation, and its std grows 11x. This same before/after pair is the evidence that R2_ens is worthless (it moved 0.997333 → 0.997341 across this collapse); see the DO-NOT-MACRO block.
% source=results/retrain_gca2_20k/symbolic/symbolic_recovery_metrics.csv filter={gCa=2.0} column=seed raw=5 n=5
\newcommand{\valRetrainRTwoCondTrueN}{5}

% note: RMSE of the neural closure itself against the true current on the fit domain. AFTER 4.059 ± 1.477 vs BEFORE 1.270 ± 0.246 — a 3.2x degradation of the learned closure. NOTE the "before" twin is NOT currently a macro; if the paper quotes the pair, add a SymbolicRmseNnTrue entry against symbolic/symbolic_recovery_metrics.csv with filter {gCa: 2.0}, the same aggregate, and `column: rmse_nn_true` (singular — see the header). Do not hand-transcribe 1.270.
% source=results/retrain_gca2_20k/symbolic/symbolic_recovery_metrics.csv filter={gCa=2.0} column=rmse_nn_true (aggregated over seed) raw={mean=4.058756617984395, std=1.476993005944329} n=5
\newcommand{\valRetrainRmseNnTrue}{4.059 \pm 1.477}

% note: RMSE of the neural closure itself against the true current on the fit domain. AFTER 4.059 ± 1.477 vs BEFORE 1.270 ± 0.246 — a 3.2x degradation of the learned closure. NOTE the "before" twin is NOT currently a macro; if the paper quotes the pair, add a SymbolicRmseNnTrue entry against symbolic/symbolic_recovery_metrics.csv with filter {gCa: 2.0}, the same aggregate, and `column: rmse_nn_true` (singular — see the header). Do not hand-transcribe 1.270.
% source=results/retrain_gca2_20k/symbolic/symbolic_recovery_metrics.csv filter={gCa=2.0} column=rmse_nn_true (aggregated over seed) raw={mean=4.058756617984395, std=1.476993005944329} n=5
\newcommand{\valRetrainRmseNnTrueMeanStd}{4.059 \pm 1.477}

% note: RMSE of the neural closure itself against the true current on the fit domain. AFTER 4.059 ± 1.477 vs BEFORE 1.270 ± 0.246 — a 3.2x degradation of the learned closure. NOTE the "before" twin is NOT currently a macro; if the paper quotes the pair, add a SymbolicRmseNnTrue entry against symbolic/symbolic_recovery_metrics.csv with filter {gCa: 2.0}, the same aggregate, and `column: rmse_nn_true` (singular — see the header). Do not hand-transcribe 1.270.
% source=results/retrain_gca2_20k/symbolic/symbolic_recovery_metrics.csv filter={gCa=2.0} column=rmse_nn_true raw=4.058756617984395 n=5
\newcommand{\valRetrainRmseNnTrueMean}{4.059}

% note: RMSE of the neural closure itself against the true current on the fit domain. AFTER 4.059 ± 1.477 vs BEFORE 1.270 ± 0.246 — a 3.2x degradation of the learned closure. NOTE the "before" twin is NOT currently a macro; if the paper quotes the pair, add a SymbolicRmseNnTrue entry against symbolic/symbolic_recovery_metrics.csv with filter {gCa: 2.0}, the same aggregate, and `column: rmse_nn_true` (singular — see the header). Do not hand-transcribe 1.270.
% source=results/retrain_gca2_20k/symbolic/symbolic_recovery_metrics.csv filter={gCa=2.0} column=rmse_nn_true raw=1.476993005944329 n=5
\newcommand{\valRetrainRmseNnTrueStd}{1.477}

% note: RMSE of the neural closure itself against the true current on the fit domain. AFTER 4.059 ± 1.477 vs BEFORE 1.270 ± 0.246 — a 3.2x degradation of the learned closure. NOTE the "before" twin is NOT currently a macro; if the paper quotes the pair, add a SymbolicRmseNnTrue entry against symbolic/symbolic_recovery_metrics.csv with filter {gCa: 2.0}, the same aggregate, and `column: rmse_nn_true` (singular — see the header). Do not hand-transcribe 1.270.
% source=results/retrain_gca2_20k/symbolic/symbolic_recovery_metrics.csv filter={gCa=2.0} column=seed raw=5 n=5
\newcommand{\valRetrainRmseNnTrueN}{5}

% note: BESPOKE SCHEMA — no filter needed (all 5 rows are gCa=2.0 full-state, one per seed) and no `window` column: the window lives in the column-name prefix. AFTER 0.756 ± 0.411 vs BEFORE UdeForecastVRmse 0.245 ± 0.028: even the trajectory forecast — the thing the loss actually optimises — got 3x worse and 15x noisier under 4x the iterations. This is the strongest single sentence against the under-training rebuttal.
% source=results/retrain_gca2_20k/metrics_retrain.csv filter={} column=forecast_V_rmse (aggregated over seed) raw={mean=0.7556217438591413, std=0.41136054939200517} n=5
\newcommand{\valRetrainForecastVRmse}{0.756 \pm 0.411}

% note: BESPOKE SCHEMA — no filter needed (all 5 rows are gCa=2.0 full-state, one per seed) and no `window` column: the window lives in the column-name prefix. AFTER 0.756 ± 0.411 vs BEFORE UdeForecastVRmse 0.245 ± 0.028: even the trajectory forecast — the thing the loss actually optimises — got 3x worse and 15x noisier under 4x the iterations. This is the strongest single sentence against the under-training rebuttal.
% source=results/retrain_gca2_20k/metrics_retrain.csv filter={} column=forecast_V_rmse (aggregated over seed) raw={mean=0.7556217438591413, std=0.41136054939200517} n=5
\newcommand{\valRetrainForecastVRmseMeanStd}{0.756 \pm 0.411}

% note: BESPOKE SCHEMA — no filter needed (all 5 rows are gCa=2.0 full-state, one per seed) and no `window` column: the window lives in the column-name prefix. AFTER 0.756 ± 0.411 vs BEFORE UdeForecastVRmse 0.245 ± 0.028: even the trajectory forecast — the thing the loss actually optimises — got 3x worse and 15x noisier under 4x the iterations. This is the strongest single sentence against the under-training rebuttal.
% source=results/retrain_gca2_20k/metrics_retrain.csv filter={} column=forecast_V_rmse raw=0.7556217438591413 n=5
\newcommand{\valRetrainForecastVRmseMean}{0.756}

% note: BESPOKE SCHEMA — no filter needed (all 5 rows are gCa=2.0 full-state, one per seed) and no `window` column: the window lives in the column-name prefix. AFTER 0.756 ± 0.411 vs BEFORE UdeForecastVRmse 0.245 ± 0.028: even the trajectory forecast — the thing the loss actually optimises — got 3x worse and 15x noisier under 4x the iterations. This is the strongest single sentence against the under-training rebuttal.
% source=results/retrain_gca2_20k/metrics_retrain.csv filter={} column=forecast_V_rmse raw=0.41136054939200517 n=5
\newcommand{\valRetrainForecastVRmseStd}{0.411}

% note: BESPOKE SCHEMA — no filter needed (all 5 rows are gCa=2.0 full-state, one per seed) and no `window` column: the window lives in the column-name prefix. AFTER 0.756 ± 0.411 vs BEFORE UdeForecastVRmse 0.245 ± 0.028: even the trajectory forecast — the thing the loss actually optimises — got 3x worse and 15x noisier under 4x the iterations. This is the strongest single sentence against the under-training rebuttal.
% source=results/retrain_gca2_20k/metrics_retrain.csv filter={} column=seed raw=5 n=5
\newcommand{\valRetrainForecastVRmseN}{5}

% note: AFTER 3.847 ± 1.465 vs BEFORE UdeForecastICaRmse 1.189 ± 0.208 — the calcium closure got 3.2x worse. Bespoke schema; no ICa_rmse_norm twin exists in this table, so the normalized before/after comparison CANNOT be made. Do not construct one by dividing by a std taken from elsewhere.
% source=results/retrain_gca2_20k/metrics_retrain.csv filter={} column=forecast_ICa_rmse (aggregated over seed) raw={mean=3.8469099071166952, std=1.464938408895774} n=5
\newcommand{\valRetrainForecastICaRmse}{3.847 \pm 1.465}

% note: AFTER 3.847 ± 1.465 vs BEFORE UdeForecastICaRmse 1.189 ± 0.208 — the calcium closure got 3.2x worse. Bespoke schema; no ICa_rmse_norm twin exists in this table, so the normalized before/after comparison CANNOT be made. Do not construct one by dividing by a std taken from elsewhere.
% source=results/retrain_gca2_20k/metrics_retrain.csv filter={} column=forecast_ICa_rmse (aggregated over seed) raw={mean=3.8469099071166952, std=1.464938408895774} n=5
\newcommand{\valRetrainForecastICaRmseMeanStd}{3.847 \pm 1.465}

% note: AFTER 3.847 ± 1.465 vs BEFORE UdeForecastICaRmse 1.189 ± 0.208 — the calcium closure got 3.2x worse. Bespoke schema; no ICa_rmse_norm twin exists in this table, so the normalized before/after comparison CANNOT be made. Do not construct one by dividing by a std taken from elsewhere.
% source=results/retrain_gca2_20k/metrics_retrain.csv filter={} column=forecast_ICa_rmse raw=3.8469099071166952 n=5
\newcommand{\valRetrainForecastICaRmseMean}{3.847}

% note: AFTER 3.847 ± 1.465 vs BEFORE UdeForecastICaRmse 1.189 ± 0.208 — the calcium closure got 3.2x worse. Bespoke schema; no ICa_rmse_norm twin exists in this table, so the normalized before/after comparison CANNOT be made. Do not construct one by dividing by a std taken from elsewhere.
% source=results/retrain_gca2_20k/metrics_retrain.csv filter={} column=forecast_ICa_rmse raw=1.464938408895774 n=5
\newcommand{\valRetrainForecastICaRmseStd}{1.465}

% note: AFTER 3.847 ± 1.465 vs BEFORE UdeForecastICaRmse 1.189 ± 0.208 — the calcium closure got 3.2x worse. Bespoke schema; no ICa_rmse_norm twin exists in this table, so the normalized before/after comparison CANNOT be made. Do not construct one by dividing by a std taken from elsewhere.
% source=results/retrain_gca2_20k/metrics_retrain.csv filter={} column=seed raw=5 n=5
\newcommand{\valRetrainForecastICaRmseN}{5}

% note: *** THE LOAD-BEARING NUMBER OF THE WHOLE ARGUMENT. *** A tight final loss (29.66 ± 2.66, 9% relative spread) plus bfgs_ok = true on all 5 seeds proves the optimiser CONVERGED. So the divergent coefficients (RetrainAHat: 1.49 ± 1.48) are not optimisation failure — five converged runs found five DIFFERENT closures that fit the same training voltage equally well. That is the definition of non-identifiability, and this macro is what turns "our numbers are noisy" into "the parameter is not determined by the data". Two caveats for the prose: (1) final_loss is SUM of squared errors on RAW coordinates, not a mean — do not compare it to any NODE loss (the NODE loss is on NORMALIZED coordinates; V contributes 10^4x less). (2) It is the last callback value over the concatenated Adam+BFGS history. bfgs_ok deserves its own value macro if the paper asserts the two-stage schedule ran — a failed BFGS polish degrades SILENTLY to the Adam result with only an @warn (train_two_stage, experiment.jl:68-97).
% source=results/retrain_gca2_20k/metrics_retrain.csv filter={} column=final_loss (aggregated over seed) raw={mean=29.65990124147882, std=2.663394792238203} n=5
\newcommand{\valRetrainFinalLoss}{29.66 \pm 2.66}

% note: *** THE LOAD-BEARING NUMBER OF THE WHOLE ARGUMENT. *** A tight final loss (29.66 ± 2.66, 9% relative spread) plus bfgs_ok = true on all 5 seeds proves the optimiser CONVERGED. So the divergent coefficients (RetrainAHat: 1.49 ± 1.48) are not optimisation failure — five converged runs found five DIFFERENT closures that fit the same training voltage equally well. That is the definition of non-identifiability, and this macro is what turns "our numbers are noisy" into "the parameter is not determined by the data". Two caveats for the prose: (1) final_loss is SUM of squared errors on RAW coordinates, not a mean — do not compare it to any NODE loss (the NODE loss is on NORMALIZED coordinates; V contributes 10^4x less). (2) It is the last callback value over the concatenated Adam+BFGS history. bfgs_ok deserves its own value macro if the paper asserts the two-stage schedule ran — a failed BFGS polish degrades SILENTLY to the Adam result with only an @warn (train_two_stage, experiment.jl:68-97).
% source=results/retrain_gca2_20k/metrics_retrain.csv filter={} column=final_loss (aggregated over seed) raw={mean=29.65990124147882, std=2.663394792238203} n=5
\newcommand{\valRetrainFinalLossMeanStd}{29.66 \pm 2.66}

% note: *** THE LOAD-BEARING NUMBER OF THE WHOLE ARGUMENT. *** A tight final loss (29.66 ± 2.66, 9% relative spread) plus bfgs_ok = true on all 5 seeds proves the optimiser CONVERGED. So the divergent coefficients (RetrainAHat: 1.49 ± 1.48) are not optimisation failure — five converged runs found five DIFFERENT closures that fit the same training voltage equally well. That is the definition of non-identifiability, and this macro is what turns "our numbers are noisy" into "the parameter is not determined by the data". Two caveats for the prose: (1) final_loss is SUM of squared errors on RAW coordinates, not a mean — do not compare it to any NODE loss (the NODE loss is on NORMALIZED coordinates; V contributes 10^4x less). (2) It is the last callback value over the concatenated Adam+BFGS history. bfgs_ok deserves its own value macro if the paper asserts the two-stage schedule ran — a failed BFGS polish degrades SILENTLY to the Adam result with only an @warn (train_two_stage, experiment.jl:68-97).
% source=results/retrain_gca2_20k/metrics_retrain.csv filter={} column=final_loss raw=29.65990124147882 n=5
\newcommand{\valRetrainFinalLossMean}{29.66}

% note: *** THE LOAD-BEARING NUMBER OF THE WHOLE ARGUMENT. *** A tight final loss (29.66 ± 2.66, 9% relative spread) plus bfgs_ok = true on all 5 seeds proves the optimiser CONVERGED. So the divergent coefficients (RetrainAHat: 1.49 ± 1.48) are not optimisation failure — five converged runs found five DIFFERENT closures that fit the same training voltage equally well. That is the definition of non-identifiability, and this macro is what turns "our numbers are noisy" into "the parameter is not determined by the data". Two caveats for the prose: (1) final_loss is SUM of squared errors on RAW coordinates, not a mean — do not compare it to any NODE loss (the NODE loss is on NORMALIZED coordinates; V contributes 10^4x less). (2) It is the last callback value over the concatenated Adam+BFGS history. bfgs_ok deserves its own value macro if the paper asserts the two-stage schedule ran — a failed BFGS polish degrades SILENTLY to the Adam result with only an @warn (train_two_stage, experiment.jl:68-97).
% source=results/retrain_gca2_20k/metrics_retrain.csv filter={} column=final_loss raw=2.663394792238203 n=5
\newcommand{\valRetrainFinalLossStd}{2.66}

% note: *** THE LOAD-BEARING NUMBER OF THE WHOLE ARGUMENT. *** A tight final loss (29.66 ± 2.66, 9% relative spread) plus bfgs_ok = true on all 5 seeds proves the optimiser CONVERGED. So the divergent coefficients (RetrainAHat: 1.49 ± 1.48) are not optimisation failure — five converged runs found five DIFFERENT closures that fit the same training voltage equally well. That is the definition of non-identifiability, and this macro is what turns "our numbers are noisy" into "the parameter is not determined by the data". Two caveats for the prose: (1) final_loss is SUM of squared errors on RAW coordinates, not a mean — do not compare it to any NODE loss (the NODE loss is on NORMALIZED coordinates; V contributes 10^4x less). (2) It is the last callback value over the concatenated Adam+BFGS history. bfgs_ok deserves its own value macro if the paper asserts the two-stage schedule ran — a failed BFGS polish degrades SILENTLY to the Adam result with only an @warn (train_two_stage, experiment.jl:68-97).
% source=results/retrain_gca2_20k/metrics_retrain.csv filter={} column=seed raw=5 n=5
\newcommand{\valRetrainFinalLossN}{5}

% note: n = 5 seeds, extracted rather than asserted, so that a dropped run cannot silently invalidate the paper's stated n. safe_run turns a failure into MISSING ROWS, not an error (experiments_runner.jl) — n_seeds is the only trace. Every "± std over N seeds" in the paper should read \valUdeSeedCount, never a literal 5. Reminder for the Methods text: a seed redraws the network init AND the OU noise realization jointly; the two variance sources are inseparable in these runs.
% source=results/metrics_baseline_multiseed.csv filter={model=UDE, observed=full, window=forecast, metric=V_rmse} column=n_seeds raw=5.0 n=1
\newcommand{\valUdeSeedCount}{5}

% note: *** THE HONESTY COUNTER. n = 2, NOT 5. *** The NODE's spike accuracy (1.272 ± 0.415 ms) is averaged over only 2 of 5 seeds: the other 3 produced no spike matchable within the 2 ms tolerance, mean_abs_dev_ms came back NaN, and agg()'s filter(isfinite, ...) silently dropped them from both the mean AND the denominator. Quoting NODE 1.27 ms beside the UDE's 0.0 ± 0.0 (n=5) without this counter would materially FLATTER a model that diverges within one sample of the forecast start. Either report n inline via this macro or drop the metric for NODE entirely; the second is cleaner. This macro exists so the choice is explicit rather than accidental.
% source=results/metrics_baseline_multiseed.csv filter={model=NODE, observed=full, window=forecast, metric=spike_mean_abs_dev_ms} column=n_seeds raw=2.0 n=1
\newcommand{\valNodeSpikeDevSeedCount}{2}

% note: The gate error of the VOLTAGE-ONLY runs — the subject of physics-prior-reconstructs-unobserved-gates. *** DO NOT substitute UdeForecastGateRmse: that macro filters observed=full, and the claim is about the runs where the gates never entered the loss. Swapping them makes the claim's own subject wrong. *** The gates are hidden from the OBJECTIVE, not from the solver (u0 is the full true 6-state IC), so say "never scored during training", never "never known".
% source=results/metrics_baseline_multiseed.csv filter={model=UDE, observed=voltage, window=forecast, metric=gate_rmse_mean} column=mean/std raw={mean=0.0010597619993722, std=0.0004335527474268237} n=5
\newcommand{\valUdeVoltageForecastGateRmse}{0.00106 \pm 0.00043}

% note: The gate error of the VOLTAGE-ONLY runs — the subject of physics-prior-reconstructs-unobserved-gates. *** DO NOT substitute UdeForecastGateRmse: that macro filters observed=full, and the claim is about the runs where the gates never entered the loss. Swapping them makes the claim's own subject wrong. *** The gates are hidden from the OBJECTIVE, not from the solver (u0 is the full true 6-state IC), so say "never scored during training", never "never known".
% source=results/metrics_baseline_multiseed.csv filter={model=UDE, observed=voltage, window=forecast, metric=gate_rmse_mean} column=mean/std raw={mean=0.0010597619993722, std=0.0004335527474268237} n=5
\newcommand{\valUdeVoltageForecastGateRmseMeanStd}{0.00106 \pm 0.00043}

% note: The gate error of the VOLTAGE-ONLY runs — the subject of physics-prior-reconstructs-unobserved-gates. *** DO NOT substitute UdeForecastGateRmse: that macro filters observed=full, and the claim is about the runs where the gates never entered the loss. Swapping them makes the claim's own subject wrong. *** The gates are hidden from the OBJECTIVE, not from the solver (u0 is the full true 6-state IC), so say "never scored during training", never "never known".
% source=results/metrics_baseline_multiseed.csv filter={model=UDE, observed=voltage, window=forecast, metric=gate_rmse_mean} column=mean raw=0.0010597619993722 n=5
\newcommand{\valUdeVoltageForecastGateRmseMean}{0.00106}

% note: The gate error of the VOLTAGE-ONLY runs — the subject of physics-prior-reconstructs-unobserved-gates. *** DO NOT substitute UdeForecastGateRmse: that macro filters observed=full, and the claim is about the runs where the gates never entered the loss. Swapping them makes the claim's own subject wrong. *** The gates are hidden from the OBJECTIVE, not from the solver (u0 is the full true 6-state IC), so say "never scored during training", never "never known".
% source=results/metrics_baseline_multiseed.csv filter={model=UDE, observed=voltage, window=forecast, metric=gate_rmse_mean} column=std raw=0.0004335527474268237 n=5
\newcommand{\valUdeVoltageForecastGateRmseStd}{0.00043}

% note: The gate error of the VOLTAGE-ONLY runs — the subject of physics-prior-reconstructs-unobserved-gates. *** DO NOT substitute UdeForecastGateRmse: that macro filters observed=full, and the claim is about the runs where the gates never entered the loss. Swapping them makes the claim's own subject wrong. *** The gates are hidden from the OBJECTIVE, not from the solver (u0 is the full true 6-state IC), so say "never scored during training", never "never known".
% source=results/metrics_baseline_multiseed.csv filter={model=UDE, observed=voltage, window=forecast, metric=gate_rmse_mean} column=n_seeds raw=5 n=5
\newcommand{\valUdeVoltageForecastGateRmseN}{5}

% note: THE CONVERGENCE CONTROL. All 5 seeds report bfgs_ok=true, which is what turns the aggressive retrain from "we under-trained" into non-identifiability: the optimizer converged, and it converged to DIFFERENT closures that fit the same voltage equally well. Pair with RetrainFinalLoss. A boolean column: emit list renders it "true, true, ..." verbatim — it is a flag, not a quantity, so never average it.
% source=results/retrain_gca2_20k/metrics_retrain.csv filter={} column=bfgs_ok raw=[True, True, True, True, True] n=5
\newcommand{\valRetrainBfgsOk}{true, true, true, true, true}

% note: Sparse regression that never sparsifies: all 6 library terms retained at every one of the 5 seeds. This is the point of sindy-does-not-sparsify — emit list, not mean_std, because "6.0 ± 0.0" would present a constant as a distribution and invite the reader to wonder about the spread. There is no spread. The 6-term library and TAU=0.05 are in src/metrics.jl.
% source=results/symbolic/symbolic_recovery_metrics.csv filter={gCa=2.0} column=n_sindy_terms raw=[6.0, 6.0, 6.0, 6.0, 6.0] n=5
\newcommand{\valSindyNTermsRetained}{6, 6, 6, 6, 6}

% note: The BEFORE half of retrain-seed-decoupling, at the representative seed 1111 (REP_SEED, experiments_runner.jl:46). The full baseline condition must be pinned — gCa/noise_level/t_train_end all vary in metrics_all.csv, and without them this filter matches 12 rows. Compare only against RetrainRepSeedForecastVRmse, which is the same seed.
% source=results/metrics_all.csv filter={model=UDE, observed=full, window=forecast, metric=V_rmse, seed=1111, gCa=2.0, noise_level=0.02, t_train_end=30.0} column=value raw=0.2191970340109573 n=1
\newcommand{\valBaselineRepSeedForecastVRmse}{0.219}

% note: The AFTER half of retrain-seed-decoupling at seed 1111. BESPOKE SCHEMA — no gCa/window columns exist here (all 5 rows are gCa=2.0 full-state, one per seed; the window lives in the column-name prefix). Paired same-seed with BaselineRepSeedForecastVRmse so the comparison is not confounded by seed. NOTE the direction: 0.219 -> 0.251 under 4x the iterations. Worse.
% source=results/retrain_gca2_20k/metrics_retrain.csv filter={seed=1111} column=forecast_V_rmse raw=0.2509357938396778 n=1
\newcommand{\valRetrainRepSeedForecastVRmse}{0.251}

% note: BEFORE, seed 1111. Quote r2_cond_true, never the ensemble R^2=0.997, which is uninformative (it moved 0.997333 -> 0.997341 while per-seed R^2 collapsed 0.984 -> 0.845). See claim-guards.md.
% source=results/symbolic/symbolic_recovery_metrics.csv filter={gCa=2.0, seed=1111} column=r2_cond_true raw=0.9881456983773574 n=1
\newcommand{\valBaselineRepSeedRTwoCondTrue}{0.988}

% note: AFTER, seed 1111: 0.988 -> 0.580. The single most legible number in the retrain: quadrupling the budget did not refine the closure, it moved to a different one. Same-seed pairing with BaselineRepSeedRTwoCondTrue is what makes this a decoupling result rather than a seed-noise result.
% source=results/retrain_gca2_20k/symbolic/symbolic_recovery_metrics.csv filter={gCa=2.0, seed=1111} column=r2_cond_true raw=0.579578157193855 n=1
\newcommand{\valRetrainRepSeedRTwoCondTrue}{0.580}

% note: The headline positive control. 2.09 +/- 0.13 against a true 2.0, a 6.2% relative spread that brackets the truth. Cross-validated three ways: the profile-likelihood 95% CI and the Gauss-Newton relative standard error (CondRelSeGcaTwo, 6.7%) both agree with this observed seed spread. Quote WITH SymbolicAHat, always.
% source=results/identifiability/parametric_fit.csv filter={gCa_true=2.0} column=gCa_hat (aggregated over seed) raw={mean=2.092384218959993, std=0.13001189990915696} n=5
\newcommand{\valParamFitGcaHatTwo}{2.092 \pm 0.130}

% note: The headline positive control. 2.09 +/- 0.13 against a true 2.0, a 6.2% relative spread that brackets the truth. Cross-validated three ways: the profile-likelihood 95% CI and the Gauss-Newton relative standard error (CondRelSeGcaTwo, 6.7%) both agree with this observed seed spread. Quote WITH SymbolicAHat, always.
% source=results/identifiability/parametric_fit.csv filter={gCa_true=2.0} column=gCa_hat (aggregated over seed) raw={mean=2.092384218959993, std=0.13001189990915696} n=5
\newcommand{\valParamFitGcaHatTwoMeanStd}{2.092 \pm 0.130}

% note: The headline positive control. 2.09 +/- 0.13 against a true 2.0, a 6.2% relative spread that brackets the truth. Cross-validated three ways: the profile-likelihood 95% CI and the Gauss-Newton relative standard error (CondRelSeGcaTwo, 6.7%) both agree with this observed seed spread. Quote WITH SymbolicAHat, always.
% source=results/identifiability/parametric_fit.csv filter={gCa_true=2.0} column=gCa_hat raw=2.092384218959993 n=5
\newcommand{\valParamFitGcaHatTwoMean}{2.092}

% note: The headline positive control. 2.09 +/- 0.13 against a true 2.0, a 6.2% relative spread that brackets the truth. Cross-validated three ways: the profile-likelihood 95% CI and the Gauss-Newton relative standard error (CondRelSeGcaTwo, 6.7%) both agree with this observed seed spread. Quote WITH SymbolicAHat, always.
% source=results/identifiability/parametric_fit.csv filter={gCa_true=2.0} column=gCa_hat raw=0.13001189990915696 n=5
\newcommand{\valParamFitGcaHatTwoStd}{0.130}

% note: The headline positive control. 2.09 +/- 0.13 against a true 2.0, a 6.2% relative spread that brackets the truth. Cross-validated three ways: the profile-likelihood 95% CI and the Gauss-Newton relative standard error (CondRelSeGcaTwo, 6.7%) both agree with this observed seed spread. Quote WITH SymbolicAHat, always.
% source=results/identifiability/parametric_fit.csv filter={gCa_true=2.0} column=seed raw=5 n=5
\newcommand{\valParamFitGcaHatTwoN}{5}

% note: The span, for direct comparison with SymbolicAHatRange (0.673--3.095). A factor of 1.18 against the closure's 4.6.
% source=results/identifiability/parametric_fit.csv filter={gCa_true=2.0} column=gCa_hat raw={min=1.8770173788332696, max=2.211209810718123} n=5
\newcommand{\valParamFitGcaHatRangeTwo}{1.877--2.211}

% note: The reversal potential is the WEAKER of the two parameters (10.7% spread vs 6.2%), consistent with the FIM relative standard errors and with the near -1 correlation between log gCa and log ECa. Do not present ECa as recovered as cleanly as gCa.
% source=results/identifiability/parametric_fit.csv filter={gCa_true=2.0} column=ECa_hat (aggregated over seed) raw={mean=111.91100463023713, std=12.039062096659686} n=5
\newcommand{\valParamFitEcaHatTwo}{111.9 \pm 12.0}

% note: The reversal potential is the WEAKER of the two parameters (10.7% spread vs 6.2%), consistent with the FIM relative standard errors and with the near -1 correlation between log gCa and log ECa. Do not present ECa as recovered as cleanly as gCa.
% source=results/identifiability/parametric_fit.csv filter={gCa_true=2.0} column=ECa_hat (aggregated over seed) raw={mean=111.91100463023713, std=12.039062096659686} n=5
\newcommand{\valParamFitEcaHatTwoMeanStd}{111.9 \pm 12.0}

% note: The reversal potential is the WEAKER of the two parameters (10.7% spread vs 6.2%), consistent with the FIM relative standard errors and with the near -1 correlation between log gCa and log ECa. Do not present ECa as recovered as cleanly as gCa.
% source=results/identifiability/parametric_fit.csv filter={gCa_true=2.0} column=ECa_hat raw=111.91100463023713 n=5
\newcommand{\valParamFitEcaHatTwoMean}{111.9}

% note: The reversal potential is the WEAKER of the two parameters (10.7% spread vs 6.2%), consistent with the FIM relative standard errors and with the near -1 correlation between log gCa and log ECa. Do not present ECa as recovered as cleanly as gCa.
% source=results/identifiability/parametric_fit.csv filter={gCa_true=2.0} column=ECa_hat raw=12.039062096659686 n=5
\newcommand{\valParamFitEcaHatTwoStd}{12.0}

% note: The reversal potential is the WEAKER of the two parameters (10.7% spread vs 6.2%), consistent with the FIM relative standard errors and with the near -1 correlation between log gCa and log ECa. Do not present ECa as recovered as cleanly as gCa.
% source=results/identifiability/parametric_fit.csv filter={gCa_true=2.0} column=seed raw=5 n=5
\newcommand{\valParamFitEcaHatTwoN}{5}

% note: Even at the PHYSIOLOGICAL conductance the parametric fit recovers gCa (0.44 +/- 0.07), where the UDE plus symbolic route returns the WRONG SIGN (SymbolicAHatNegControl = -0.37). The contrast at gCa=0.4 is starker than at 2.0 and is the sharpest single statement of the paper's thesis.
% source=results/identifiability/parametric_fit.csv filter={gCa_true=0.4} column=gCa_hat (aggregated over seed) raw={mean=0.43775053086048327, std=0.06510726518191451} n=5
\newcommand{\valParamFitGcaHatPointFour}{0.438 \pm 0.065}

% note: Even at the PHYSIOLOGICAL conductance the parametric fit recovers gCa (0.44 +/- 0.07), where the UDE plus symbolic route returns the WRONG SIGN (SymbolicAHatNegControl = -0.37). The contrast at gCa=0.4 is starker than at 2.0 and is the sharpest single statement of the paper's thesis.
% source=results/identifiability/parametric_fit.csv filter={gCa_true=0.4} column=gCa_hat (aggregated over seed) raw={mean=0.43775053086048327, std=0.06510726518191451} n=5
\newcommand{\valParamFitGcaHatPointFourMeanStd}{0.438 \pm 0.065}

% note: Even at the PHYSIOLOGICAL conductance the parametric fit recovers gCa (0.44 +/- 0.07), where the UDE plus symbolic route returns the WRONG SIGN (SymbolicAHatNegControl = -0.37). The contrast at gCa=0.4 is starker than at 2.0 and is the sharpest single statement of the paper's thesis.
% source=results/identifiability/parametric_fit.csv filter={gCa_true=0.4} column=gCa_hat raw=0.43775053086048327 n=5
\newcommand{\valParamFitGcaHatPointFourMean}{0.438}

% note: Even at the PHYSIOLOGICAL conductance the parametric fit recovers gCa (0.44 +/- 0.07), where the UDE plus symbolic route returns the WRONG SIGN (SymbolicAHatNegControl = -0.37). The contrast at gCa=0.4 is starker than at 2.0 and is the sharpest single statement of the paper's thesis.
% source=results/identifiability/parametric_fit.csv filter={gCa_true=0.4} column=gCa_hat raw=0.06510726518191451 n=5
\newcommand{\valParamFitGcaHatPointFourStd}{0.065}

% note: Even at the PHYSIOLOGICAL conductance the parametric fit recovers gCa (0.44 +/- 0.07), where the UDE plus symbolic route returns the WRONG SIGN (SymbolicAHatNegControl = -0.37). The contrast at gCa=0.4 is starker than at 2.0 and is the sharpest single statement of the paper's thesis.
% source=results/identifiability/parametric_fit.csv filter={gCa_true=0.4} column=seed raw=5 n=5
\newcommand{\valParamFitGcaHatPointFourN}{5}

% source=results/identifiability/parametric_fit.csv filter={gCa_true=1.0} column=gCa_hat (aggregated over seed) raw={mean=1.0430471597691386, std=0.0840708259850014} n=5
\newcommand{\valParamFitGcaHatOne}{1.043 \pm 0.084}

% source=results/identifiability/parametric_fit.csv filter={gCa_true=1.0} column=gCa_hat (aggregated over seed) raw={mean=1.0430471597691386, std=0.0840708259850014} n=5
\newcommand{\valParamFitGcaHatOneMeanStd}{1.043 \pm 0.084}

% source=results/identifiability/parametric_fit.csv filter={gCa_true=1.0} column=gCa_hat raw=1.0430471597691386 n=5
\newcommand{\valParamFitGcaHatOneMean}{1.043}

% source=results/identifiability/parametric_fit.csv filter={gCa_true=1.0} column=gCa_hat raw=0.0840708259850014 n=5
\newcommand{\valParamFitGcaHatOneStd}{0.084}

% source=results/identifiability/parametric_fit.csv filter={gCa_true=1.0} column=seed raw=5 n=5
\newcommand{\valParamFitGcaHatOneN}{5}

% note: Completes the monotone tightening across the sweep: relative spread runs 14.9% then 8.1% then 6.2% then 2.9% as the calcium signal grows. The parametric fit degrades gracefully where the closure degrades catastrophically.
% source=results/identifiability/parametric_fit.csv filter={gCa_true=4.0} column=gCa_hat (aggregated over seed) raw={mean=4.10529980106318, std=0.11984914182199823} n=5
\newcommand{\valParamFitGcaHatFour}{4.105 \pm 0.120}

% note: Completes the monotone tightening across the sweep: relative spread runs 14.9% then 8.1% then 6.2% then 2.9% as the calcium signal grows. The parametric fit degrades gracefully where the closure degrades catastrophically.
% source=results/identifiability/parametric_fit.csv filter={gCa_true=4.0} column=gCa_hat (aggregated over seed) raw={mean=4.10529980106318, std=0.11984914182199823} n=5
\newcommand{\valParamFitGcaHatFourMeanStd}{4.105 \pm 0.120}

% note: Completes the monotone tightening across the sweep: relative spread runs 14.9% then 8.1% then 6.2% then 2.9% as the calcium signal grows. The parametric fit degrades gracefully where the closure degrades catastrophically.
% source=results/identifiability/parametric_fit.csv filter={gCa_true=4.0} column=gCa_hat raw=4.10529980106318 n=5
\newcommand{\valParamFitGcaHatFourMean}{4.105}

% note: Completes the monotone tightening across the sweep: relative spread runs 14.9% then 8.1% then 6.2% then 2.9% as the calcium signal grows. The parametric fit degrades gracefully where the closure degrades catastrophically.
% source=results/identifiability/parametric_fit.csv filter={gCa_true=4.0} column=gCa_hat raw=0.11984914182199823 n=5
\newcommand{\valParamFitGcaHatFourStd}{0.120}

% note: Completes the monotone tightening across the sweep: relative spread runs 14.9% then 8.1% then 6.2% then 2.9% as the calcium signal grows. The parametric fit degrades gracefully where the closure degrades catastrophically.
% source=results/identifiability/parametric_fit.csv filter={gCa_true=4.0} column=seed raw=5 n=5
\newcommand{\valParamFitGcaHatFourN}{5}

% note: chi-square per data point at the optimum, with sigma_i the KNOWN per-channel noise SD. A value of 1.0 means the fit sits exactly at the noise floor: the model is correct and the weighting is right. This is what licenses the chi-square threshold used for the profile-likelihood CI. Report it as a goodness-of-fit sanity check, not as a result in itself.
% source=results/identifiability/parametric_fit.csv filter={gCa_true=2.0} column=chi2_per_point raw=0.9944792850770685 n=5
\newcommand{\valParamFitChiSqPerPointTwo}{0.994}

% note: The parametric fit ALSO forecasts about twice as well as the UDE. Quote this whenever the comparison is made: it is the reason the honest claim is "the data determine the parameter" and NOT "two equally good fits disagree". Omitting it would overstate the result.
% source=results/identifiability/parametric_fit.csv filter={gCa_true=2.0} column=forecast_V_rmse (aggregated over seed) raw={mean=0.12296809823224704, std=0.031883617114038566} n=5
\newcommand{\valParamFitForecastVRmseTwo}{0.123 \pm 0.032}

% note: The parametric fit ALSO forecasts about twice as well as the UDE. Quote this whenever the comparison is made: it is the reason the honest claim is "the data determine the parameter" and NOT "two equally good fits disagree". Omitting it would overstate the result.
% source=results/identifiability/parametric_fit.csv filter={gCa_true=2.0} column=forecast_V_rmse (aggregated over seed) raw={mean=0.12296809823224704, std=0.031883617114038566} n=5
\newcommand{\valParamFitForecastVRmseTwoMeanStd}{0.123 \pm 0.032}

% note: The parametric fit ALSO forecasts about twice as well as the UDE. Quote this whenever the comparison is made: it is the reason the honest claim is "the data determine the parameter" and NOT "two equally good fits disagree". Omitting it would overstate the result.
% source=results/identifiability/parametric_fit.csv filter={gCa_true=2.0} column=forecast_V_rmse raw=0.12296809823224704 n=5
\newcommand{\valParamFitForecastVRmseTwoMean}{0.123}

% note: The parametric fit ALSO forecasts about twice as well as the UDE. Quote this whenever the comparison is made: it is the reason the honest claim is "the data determine the parameter" and NOT "two equally good fits disagree". Omitting it would overstate the result.
% source=results/identifiability/parametric_fit.csv filter={gCa_true=2.0} column=forecast_V_rmse raw=0.031883617114038566 n=5
\newcommand{\valParamFitForecastVRmseTwoStd}{0.032}

% note: The parametric fit ALSO forecasts about twice as well as the UDE. Quote this whenever the comparison is made: it is the reason the honest claim is "the data determine the parameter" and NOT "two equally good fits disagree". Omitting it would overstate the result.
% source=results/identifiability/parametric_fit.csv filter={gCa_true=2.0} column=seed raw=5 n=5
\newcommand{\valParamFitForecastVRmseTwoN}{5}

% note: Asymptotic RELATIVE standard error of gCa, 6.7%. Matches the observed seed spread (6.2%, ParamFitGcaHatTwo) almost exactly, an independent confirmation that the seed-to-seed scatter of the parametric fit is ordinary estimation noise and nothing more.
% source=results/identifiability/conditioning.csv filter={gCa_true=2.0} column=rel_se_gCa (aggregated over seed) raw={mean=0.06667186168953679, std=0.004331278030106411} n=5
\newcommand{\valCondRelSeGcaTwo}{0.067 \pm 0.004}

% note: Asymptotic RELATIVE standard error of gCa, 6.7%. Matches the observed seed spread (6.2%, ParamFitGcaHatTwo) almost exactly, an independent confirmation that the seed-to-seed scatter of the parametric fit is ordinary estimation noise and nothing more.
% source=results/identifiability/conditioning.csv filter={gCa_true=2.0} column=rel_se_gCa (aggregated over seed) raw={mean=0.06667186168953679, std=0.004331278030106411} n=5
\newcommand{\valCondRelSeGcaTwoMeanStd}{0.067 \pm 0.004}

% note: Asymptotic RELATIVE standard error of gCa, 6.7%. Matches the observed seed spread (6.2%, ParamFitGcaHatTwo) almost exactly, an independent confirmation that the seed-to-seed scatter of the parametric fit is ordinary estimation noise and nothing more.
% source=results/identifiability/conditioning.csv filter={gCa_true=2.0} column=rel_se_gCa raw=0.06667186168953679 n=5
\newcommand{\valCondRelSeGcaTwoMean}{0.067}

% note: Asymptotic RELATIVE standard error of gCa, 6.7%. Matches the observed seed spread (6.2%, ParamFitGcaHatTwo) almost exactly, an independent confirmation that the seed-to-seed scatter of the parametric fit is ordinary estimation noise and nothing more.
% source=results/identifiability/conditioning.csv filter={gCa_true=2.0} column=rel_se_gCa raw=0.004331278030106411 n=5
\newcommand{\valCondRelSeGcaTwoStd}{0.004}

% note: Asymptotic RELATIVE standard error of gCa, 6.7%. Matches the observed seed spread (6.2%, ParamFitGcaHatTwo) almost exactly, an independent confirmation that the seed-to-seed scatter of the parametric fit is ordinary estimation noise and nothing more.
% source=results/identifiability/conditioning.csv filter={gCa_true=2.0} column=seed raw=5 n=5
\newcommand{\valCondRelSeGcaTwoN}{5}

% source=results/identifiability/conditioning.csv filter={gCa_true=2.0} column=rel_se_ECa (aggregated over seed) raw={mean=0.11104759082026801, std=0.002608339618740706} n=5
\newcommand{\valCondRelSeEcaTwo}{0.111 \pm 0.003}

% source=results/identifiability/conditioning.csv filter={gCa_true=2.0} column=rel_se_ECa (aggregated over seed) raw={mean=0.11104759082026801, std=0.002608339618740706} n=5
\newcommand{\valCondRelSeEcaTwoMeanStd}{0.111 \pm 0.003}

% source=results/identifiability/conditioning.csv filter={gCa_true=2.0} column=rel_se_ECa raw=0.11104759082026801 n=5
\newcommand{\valCondRelSeEcaTwoMean}{0.111}

% source=results/identifiability/conditioning.csv filter={gCa_true=2.0} column=rel_se_ECa raw=0.002608339618740706 n=5
\newcommand{\valCondRelSeEcaTwoStd}{0.003}

% source=results/identifiability/conditioning.csv filter={gCa_true=2.0} column=seed raw=5 n=5
\newcommand{\valCondRelSeEcaTwoN}{5}

% note: At the physiological conductance the parametric fit is four times less precise (26.9% vs 6.7%) but still finite and still correctly signed. The information degrades smoothly with signal amplitude; it does not vanish.
% source=results/identifiability/conditioning.csv filter={gCa_true=0.4} column=rel_se_gCa (aggregated over seed) raw={mean=0.2694265739421785, std=0.042432531341401535} n=5
\newcommand{\valCondRelSeGcaPointFour}{0.269 \pm 0.042}

% note: At the physiological conductance the parametric fit is four times less precise (26.9% vs 6.7%) but still finite and still correctly signed. The information degrades smoothly with signal amplitude; it does not vanish.
% source=results/identifiability/conditioning.csv filter={gCa_true=0.4} column=rel_se_gCa (aggregated over seed) raw={mean=0.2694265739421785, std=0.042432531341401535} n=5
\newcommand{\valCondRelSeGcaPointFourMeanStd}{0.269 \pm 0.042}

% note: At the physiological conductance the parametric fit is four times less precise (26.9% vs 6.7%) but still finite and still correctly signed. The information degrades smoothly with signal amplitude; it does not vanish.
% source=results/identifiability/conditioning.csv filter={gCa_true=0.4} column=rel_se_gCa raw=0.2694265739421785 n=5
\newcommand{\valCondRelSeGcaPointFourMean}{0.269}

% note: At the physiological conductance the parametric fit is four times less precise (26.9% vs 6.7%) but still finite and still correctly signed. The information degrades smoothly with signal amplitude; it does not vanish.
% source=results/identifiability/conditioning.csv filter={gCa_true=0.4} column=rel_se_gCa raw=0.042432531341401535 n=5
\newcommand{\valCondRelSeGcaPointFourStd}{0.042}

% note: At the physiological conductance the parametric fit is four times less precise (26.9% vs 6.7%) but still finite and still correctly signed. The information degrades smoothly with signal amplitude; it does not vanish.
% source=results/identifiability/conditioning.csv filter={gCa_true=0.4} column=seed raw=5 n=5
\newcommand{\valCondRelSeGcaPointFourN}{5}

% note: Condition number of the log-parameter Fisher information. Large, so there IS a sloppy direction: gCa and ECa trade off along an identifiable combination (their log-correlation is about -1.00). But sloppy is not the same as non-identifiable, since both parameters still have finite, small relative standard errors. That distinction is the point.
% source=results/identifiability/conditioning.csv filter={gCa_true=2.0} column=cond_number raw=7087.976234532391 n=5
\newcommand{\valCondNumberTwo}{7088}

% note: Fills the "not recorded" cell in the retrain table. Directly comparable to TrainLossAfter: identical objective, data, seeds, solver and tolerances, with only the iteration budget differing.
% source=results/identifiability/training_losses_before_after.csv filter={} column=final_loss_before (aggregated over seed) raw={mean=35.71475238532089, std=1.8135366063932328} n=5
\newcommand{\valTrainLossBefore}{35.71 \pm 1.81}

% note: Fills the "not recorded" cell in the retrain table. Directly comparable to TrainLossAfter: identical objective, data, seeds, solver and tolerances, with only the iteration budget differing.
% source=results/identifiability/training_losses_before_after.csv filter={} column=final_loss_before (aggregated over seed) raw={mean=35.71475238532089, std=1.8135366063932328} n=5
\newcommand{\valTrainLossBeforeMeanStd}{35.71 \pm 1.81}

% note: Fills the "not recorded" cell in the retrain table. Directly comparable to TrainLossAfter: identical objective, data, seeds, solver and tolerances, with only the iteration budget differing.
% source=results/identifiability/training_losses_before_after.csv filter={} column=final_loss_before raw=35.71475238532089 n=5
\newcommand{\valTrainLossBeforeMean}{35.71}

% note: Fills the "not recorded" cell in the retrain table. Directly comparable to TrainLossAfter: identical objective, data, seeds, solver and tolerances, with only the iteration budget differing.
% source=results/identifiability/training_losses_before_after.csv filter={} column=final_loss_before raw=1.8135366063932328 n=5
\newcommand{\valTrainLossBeforeStd}{1.81}

% note: Fills the "not recorded" cell in the retrain table. Directly comparable to TrainLossAfter: identical objective, data, seeds, solver and tolerances, with only the iteration budget differing.
% source=results/identifiability/training_losses_before_after.csv filter={} column=seed raw=5 n=5
\newcommand{\valTrainLossBeforeN}{5}

% note: *** THE RETRAIN GENUINELY OPTIMISED BETTER: 35.71 to 29.66, a 17% reduction. *** This STRENGTHENS the control rather than weakening it. The earlier framing ("the final losses are tight, so the optimiser converged") understated the result: the aggressive run demonstrably fit the training data better while the conductance estimate, the R-squared and the forecast all got worse. That is optimisation along a flat direction, seen directly.
% source=results/identifiability/training_losses_before_after.csv filter={} column=final_loss_after_recomputed (aggregated over seed) raw={mean=29.65990124147882, std=2.663394792238203} n=5
\newcommand{\valTrainLossAfter}{29.66 \pm 2.66}

% note: *** THE RETRAIN GENUINELY OPTIMISED BETTER: 35.71 to 29.66, a 17% reduction. *** This STRENGTHENS the control rather than weakening it. The earlier framing ("the final losses are tight, so the optimiser converged") understated the result: the aggressive run demonstrably fit the training data better while the conductance estimate, the R-squared and the forecast all got worse. That is optimisation along a flat direction, seen directly.
% source=results/identifiability/training_losses_before_after.csv filter={} column=final_loss_after_recomputed (aggregated over seed) raw={mean=29.65990124147882, std=2.663394792238203} n=5
\newcommand{\valTrainLossAfterMeanStd}{29.66 \pm 2.66}

% note: *** THE RETRAIN GENUINELY OPTIMISED BETTER: 35.71 to 29.66, a 17% reduction. *** This STRENGTHENS the control rather than weakening it. The earlier framing ("the final losses are tight, so the optimiser converged") understated the result: the aggressive run demonstrably fit the training data better while the conductance estimate, the R-squared and the forecast all got worse. That is optimisation along a flat direction, seen directly.
% source=results/identifiability/training_losses_before_after.csv filter={} column=final_loss_after_recomputed raw=29.65990124147882 n=5
\newcommand{\valTrainLossAfterMean}{29.66}

% note: *** THE RETRAIN GENUINELY OPTIMISED BETTER: 35.71 to 29.66, a 17% reduction. *** This STRENGTHENS the control rather than weakening it. The earlier framing ("the final losses are tight, so the optimiser converged") understated the result: the aggressive run demonstrably fit the training data better while the conductance estimate, the R-squared and the forecast all got worse. That is optimisation along a flat direction, seen directly.
% source=results/identifiability/training_losses_before_after.csv filter={} column=final_loss_after_recomputed raw=2.663394792238203 n=5
\newcommand{\valTrainLossAfterStd}{2.66}

% note: *** THE RETRAIN GENUINELY OPTIMISED BETTER: 35.71 to 29.66, a 17% reduction. *** This STRENGTHENS the control rather than weakening it. The earlier framing ("the final losses are tight, so the optimiser converged") understated the result: the aggressive run demonstrably fit the training data better while the conductance estimate, the R-squared and the forecast all got worse. That is optimisation along a flat direction, seen directly.
% source=results/identifiability/training_losses_before_after.csv filter={} column=seed raw=5 n=5
\newcommand{\valTrainLossAfterN}{5}

% note: *** PREFER THIS OVER SymbolicAHat AS THE HEADLINE. *** Same estimator, but fitted only on points the model was actually supervised on (t <= 30 ms, n=154). It concedes the hull objection completely and costs nothing: the estimate is if anything slightly closer to the truth than the hull fit, and the spread is comparable. Keep the hull only for the SINDy comparison, where its conditioning rationale genuinely applies.
% source=results/identifiability/symbolic_domain_comparison.csv filter={domain=traj-train} column=a_hat (aggregated over seed) raw={mean=1.6445778751619045, std=0.6717585819876875} n=5
\newcommand{\valSymbolicAHatTrajTrain}{1.645 \pm 0.672}

% note: *** PREFER THIS OVER SymbolicAHat AS THE HEADLINE. *** Same estimator, but fitted only on points the model was actually supervised on (t <= 30 ms, n=154). It concedes the hull objection completely and costs nothing: the estimate is if anything slightly closer to the truth than the hull fit, and the spread is comparable. Keep the hull only for the SINDy comparison, where its conditioning rationale genuinely applies.
% source=results/identifiability/symbolic_domain_comparison.csv filter={domain=traj-train} column=a_hat (aggregated over seed) raw={mean=1.6445778751619045, std=0.6717585819876875} n=5
\newcommand{\valSymbolicAHatTrajTrainMeanStd}{1.645 \pm 0.672}

% note: *** PREFER THIS OVER SymbolicAHat AS THE HEADLINE. *** Same estimator, but fitted only on points the model was actually supervised on (t <= 30 ms, n=154). It concedes the hull objection completely and costs nothing: the estimate is if anything slightly closer to the truth than the hull fit, and the spread is comparable. Keep the hull only for the SINDy comparison, where its conditioning rationale genuinely applies.
% source=results/identifiability/symbolic_domain_comparison.csv filter={domain=traj-train} column=a_hat raw=1.6445778751619045 n=5
\newcommand{\valSymbolicAHatTrajTrainMean}{1.645}

% note: *** PREFER THIS OVER SymbolicAHat AS THE HEADLINE. *** Same estimator, but fitted only on points the model was actually supervised on (t <= 30 ms, n=154). It concedes the hull objection completely and costs nothing: the estimate is if anything slightly closer to the truth than the hull fit, and the spread is comparable. Keep the hull only for the SINDy comparison, where its conditioning rationale genuinely applies.
% source=results/identifiability/symbolic_domain_comparison.csv filter={domain=traj-train} column=a_hat raw=0.6717585819876875 n=5
\newcommand{\valSymbolicAHatTrajTrainStd}{0.672}

% note: *** PREFER THIS OVER SymbolicAHat AS THE HEADLINE. *** Same estimator, but fitted only on points the model was actually supervised on (t <= 30 ms, n=154). It concedes the hull objection completely and costs nothing: the estimate is if anything slightly closer to the truth than the hull fit, and the spread is comparable. Keep the hull only for the SINDy comparison, where its conditioning rationale genuinely applies.
% source=results/identifiability/symbolic_domain_comparison.csv filter={domain=traj-train} column=seed raw=5 n=5
\newcommand{\valSymbolicAHatTrajTrainN}{5}

% source=results/identifiability/symbolic_domain_comparison.csv filter={domain=traj-train} column=a_hat raw={min=1.2408727938278845, max=2.8156936899025493} n=5
\newcommand{\valSymbolicAHatRangeTrajTrain}{1.241--2.816}

% note: The fitted CURVE matches the true current well on the supervised trajectory (R-squared 0.99) while the COEFFICIENT that curve implies is off by 18% with a 41% spread. Curve-fit quality and coefficient identifiability are different properties; this pair of macros is the cleanest place in the paper to say so.
% source=results/identifiability/symbolic_domain_comparison.csv filter={domain=traj-train} column=r2_vs_true_on_domain (aggregated over seed) raw={mean=0.9905859909037261, std=0.0045213035318677555} n=5
\newcommand{\valSymbolicRTwoTrajTrain}{0.991 \pm 0.005}

% note: The fitted CURVE matches the true current well on the supervised trajectory (R-squared 0.99) while the COEFFICIENT that curve implies is off by 18% with a 41% spread. Curve-fit quality and coefficient identifiability are different properties; this pair of macros is the cleanest place in the paper to say so.
% source=results/identifiability/symbolic_domain_comparison.csv filter={domain=traj-train} column=r2_vs_true_on_domain (aggregated over seed) raw={mean=0.9905859909037261, std=0.0045213035318677555} n=5
\newcommand{\valSymbolicRTwoTrajTrainMeanStd}{0.991 \pm 0.005}

% note: The fitted CURVE matches the true current well on the supervised trajectory (R-squared 0.99) while the COEFFICIENT that curve implies is off by 18% with a 41% spread. Curve-fit quality and coefficient identifiability are different properties; this pair of macros is the cleanest place in the paper to say so.
% source=results/identifiability/symbolic_domain_comparison.csv filter={domain=traj-train} column=r2_vs_true_on_domain raw=0.9905859909037261 n=5
\newcommand{\valSymbolicRTwoTrajTrainMean}{0.991}

% note: The fitted CURVE matches the true current well on the supervised trajectory (R-squared 0.99) while the COEFFICIENT that curve implies is off by 18% with a 41% spread. Curve-fit quality and coefficient identifiability are different properties; this pair of macros is the cleanest place in the paper to say so.
% source=results/identifiability/symbolic_domain_comparison.csv filter={domain=traj-train} column=r2_vs_true_on_domain raw=0.0045213035318677555 n=5
\newcommand{\valSymbolicRTwoTrajTrainStd}{0.005}

% note: The fitted CURVE matches the true current well on the supervised trajectory (R-squared 0.99) while the COEFFICIENT that curve implies is off by 18% with a 41% spread. Curve-fit quality and coefficient identifiability are different properties; this pair of macros is the cleanest place in the paper to say so.
% source=results/identifiability/symbolic_domain_comparison.csv filter={domain=traj-train} column=seed raw=5 n=5
\newcommand{\valSymbolicRTwoTrajTrainN}{5}

% note: Corrects a mis-stated justification in the earlier draft. For the TWO-term conductance form the trajectory is perfectly well conditioned (about 102, against 70 on the hull), so the near-collinearity argument does NOT apply here. It applies to the SIX-term SINDy library, which is ill-conditioned on every domain (about 25k on the trajectory, 14k on the hull).
% source=results/identifiability/symbolic_domain_comparison.csv filter={domain=traj-train} column=cond raw=102.43602726682067 n=5
\newcommand{\valSymbolicCondTrajTrain}{102.4}

% source=results/identifiability/symbolic_domain_comparison.csv filter={domain=hull} column=cond raw=69.86974830325174 n=5
\newcommand{\valSymbolicCondHull}{69.9}

% note: Replaces the single-seed 1.179. The seed spread here (+/-0.318) is 13x the spread at the gCa=2.0 centre point (+/-0.025), so the earlier practice of judging every sweep point against the centre-point spread understated the uncertainty at the low end badly.
% source=results/gca_sweep_5seed/gca_sweep_multiseed.csv filter={gCa=0.4, metric=ICa_rmse_norm} column=mean/sd raw={mean=0.7242919047642109, std=0.3184350839738513} n=5
\newcommand{\valGcaSweepICaNormPointFour}{0.724 \pm 0.318}

% note: Replaces the single-seed 1.179. The seed spread here (+/-0.318) is 13x the spread at the gCa=2.0 centre point (+/-0.025), so the earlier practice of judging every sweep point against the centre-point spread understated the uncertainty at the low end badly.
% source=results/gca_sweep_5seed/gca_sweep_multiseed.csv filter={gCa=0.4, metric=ICa_rmse_norm} column=mean/sd raw={mean=0.7242919047642109, std=0.3184350839738513} n=5
\newcommand{\valGcaSweepICaNormPointFourMeanStd}{0.724 \pm 0.318}

% note: Replaces the single-seed 1.179. The seed spread here (+/-0.318) is 13x the spread at the gCa=2.0 centre point (+/-0.025), so the earlier practice of judging every sweep point against the centre-point spread understated the uncertainty at the low end badly.
% source=results/gca_sweep_5seed/gca_sweep_multiseed.csv filter={gCa=0.4, metric=ICa_rmse_norm} column=mean raw=0.7242919047642109 n=5
\newcommand{\valGcaSweepICaNormPointFourMean}{0.724}

% note: Replaces the single-seed 1.179. The seed spread here (+/-0.318) is 13x the spread at the gCa=2.0 centre point (+/-0.025), so the earlier practice of judging every sweep point against the centre-point spread understated the uncertainty at the low end badly.
% source=results/gca_sweep_5seed/gca_sweep_multiseed.csv filter={gCa=0.4, metric=ICa_rmse_norm} column=sd raw=0.3184350839738513 n=5
\newcommand{\valGcaSweepICaNormPointFourStd}{0.318}

% note: Replaces the single-seed 1.179. The seed spread here (+/-0.318) is 13x the spread at the gCa=2.0 centre point (+/-0.025), so the earlier practice of judging every sweep point against the centre-point spread understated the uncertainty at the low end badly.
% source=results/gca_sweep_5seed/gca_sweep_multiseed.csv filter={gCa=0.4, metric=ICa_rmse_norm} column=n raw=5 n=5
\newcommand{\valGcaSweepICaNormPointFourN}{5}

% source=results/gca_sweep_5seed/gca_sweep_multiseed.csv filter={gCa=1.0, metric=ICa_rmse_norm} column=mean/sd raw={mean=0.31196923397866433, std=0.09658260510624679} n=5
\newcommand{\valGcaSweepICaNormOne}{0.312 \pm 0.097}

% source=results/gca_sweep_5seed/gca_sweep_multiseed.csv filter={gCa=1.0, metric=ICa_rmse_norm} column=mean/sd raw={mean=0.31196923397866433, std=0.09658260510624679} n=5
\newcommand{\valGcaSweepICaNormOneMeanStd}{0.312 \pm 0.097}

% source=results/gca_sweep_5seed/gca_sweep_multiseed.csv filter={gCa=1.0, metric=ICa_rmse_norm} column=mean raw=0.31196923397866433 n=5
\newcommand{\valGcaSweepICaNormOneMean}{0.312}

% source=results/gca_sweep_5seed/gca_sweep_multiseed.csv filter={gCa=1.0, metric=ICa_rmse_norm} column=sd raw=0.09658260510624679 n=5
\newcommand{\valGcaSweepICaNormOneStd}{0.097}

% source=results/gca_sweep_5seed/gca_sweep_multiseed.csv filter={gCa=1.0, metric=ICa_rmse_norm} column=n raw=5 n=5
\newcommand{\valGcaSweepICaNormOneN}{5}

% source=results/gca_sweep_5seed/gca_sweep_multiseed.csv filter={gCa=2.0, metric=ICa_rmse_norm} column=mean/sd raw={mean=0.14034112152546374, std=0.024583934920143758} n=5
\newcommand{\valGcaSweepICaNormTwo}{0.140 \pm 0.025}

% source=results/gca_sweep_5seed/gca_sweep_multiseed.csv filter={gCa=2.0, metric=ICa_rmse_norm} column=mean/sd raw={mean=0.14034112152546374, std=0.024583934920143758} n=5
\newcommand{\valGcaSweepICaNormTwoMeanStd}{0.140 \pm 0.025}

% source=results/gca_sweep_5seed/gca_sweep_multiseed.csv filter={gCa=2.0, metric=ICa_rmse_norm} column=mean raw=0.14034112152546374 n=5
\newcommand{\valGcaSweepICaNormTwoMean}{0.140}

% source=results/gca_sweep_5seed/gca_sweep_multiseed.csv filter={gCa=2.0, metric=ICa_rmse_norm} column=sd raw=0.024583934920143758 n=5
\newcommand{\valGcaSweepICaNormTwoStd}{0.025}

% source=results/gca_sweep_5seed/gca_sweep_multiseed.csv filter={gCa=2.0, metric=ICa_rmse_norm} column=n raw=5 n=5
\newcommand{\valGcaSweepICaNormTwoN}{5}

% source=results/gca_sweep_5seed/gca_sweep_multiseed.csv filter={gCa=4.0, metric=ICa_rmse_norm} column=mean/sd raw={mean=0.09324722768825978, std=0.0441092449693115} n=5
\newcommand{\valGcaSweepICaNormFour}{0.093 \pm 0.044}

% source=results/gca_sweep_5seed/gca_sweep_multiseed.csv filter={gCa=4.0, metric=ICa_rmse_norm} column=mean/sd raw={mean=0.09324722768825978, std=0.0441092449693115} n=5
\newcommand{\valGcaSweepICaNormFourMeanStd}{0.093 \pm 0.044}

% source=results/gca_sweep_5seed/gca_sweep_multiseed.csv filter={gCa=4.0, metric=ICa_rmse_norm} column=mean raw=0.09324722768825978 n=5
\newcommand{\valGcaSweepICaNormFourMean}{0.093}

% source=results/gca_sweep_5seed/gca_sweep_multiseed.csv filter={gCa=4.0, metric=ICa_rmse_norm} column=sd raw=0.0441092449693115 n=5
\newcommand{\valGcaSweepICaNormFourStd}{0.044}

% source=results/gca_sweep_5seed/gca_sweep_multiseed.csv filter={gCa=4.0, metric=ICa_rmse_norm} column=n raw=5 n=5
\newcommand{\valGcaSweepICaNormFourN}{5}

% note: Pair with GcaSweepICaRawPointFour. The RAW calcium error is statistically flat across the whole sweep (1.19 to 1.48, every interval overlapping): the closure's absolute error barely changes and only the signal it must explain grows. That is a cleaner statement of the identifiability axis than the normalised metric alone, and it only became visible with 5 seeds.
% source=results/gca_sweep_5seed/gca_sweep_multiseed.csv filter={gCa=2.0, metric=ICa_rmse} column=mean/sd raw={mean=1.1890720989803751, std=0.2082929848283161} n=5
\newcommand{\valGcaSweepICaRawTwo}{1.189 \pm 0.208}

% note: Pair with GcaSweepICaRawPointFour. The RAW calcium error is statistically flat across the whole sweep (1.19 to 1.48, every interval overlapping): the closure's absolute error barely changes and only the signal it must explain grows. That is a cleaner statement of the identifiability axis than the normalised metric alone, and it only became visible with 5 seeds.
% source=results/gca_sweep_5seed/gca_sweep_multiseed.csv filter={gCa=2.0, metric=ICa_rmse} column=mean/sd raw={mean=1.1890720989803751, std=0.2082929848283161} n=5
\newcommand{\valGcaSweepICaRawTwoMeanStd}{1.189 \pm 0.208}

% note: Pair with GcaSweepICaRawPointFour. The RAW calcium error is statistically flat across the whole sweep (1.19 to 1.48, every interval overlapping): the closure's absolute error barely changes and only the signal it must explain grows. That is a cleaner statement of the identifiability axis than the normalised metric alone, and it only became visible with 5 seeds.
% source=results/gca_sweep_5seed/gca_sweep_multiseed.csv filter={gCa=2.0, metric=ICa_rmse} column=mean raw=1.1890720989803751 n=5
\newcommand{\valGcaSweepICaRawTwoMean}{1.189}

% note: Pair with GcaSweepICaRawPointFour. The RAW calcium error is statistically flat across the whole sweep (1.19 to 1.48, every interval overlapping): the closure's absolute error barely changes and only the signal it must explain grows. That is a cleaner statement of the identifiability axis than the normalised metric alone, and it only became visible with 5 seeds.
% source=results/gca_sweep_5seed/gca_sweep_multiseed.csv filter={gCa=2.0, metric=ICa_rmse} column=sd raw=0.2082929848283161 n=5
\newcommand{\valGcaSweepICaRawTwoStd}{0.208}

% note: Pair with GcaSweepICaRawPointFour. The RAW calcium error is statistically flat across the whole sweep (1.19 to 1.48, every interval overlapping): the closure's absolute error barely changes and only the signal it must explain grows. That is a cleaner statement of the identifiability axis than the normalised metric alone, and it only became visible with 5 seeds.
% source=results/gca_sweep_5seed/gca_sweep_multiseed.csv filter={gCa=2.0, metric=ICa_rmse} column=n raw=5 n=5
\newcommand{\valGcaSweepICaRawTwoN}{5}

% source=results/gca_sweep_5seed/gca_sweep_multiseed.csv filter={gCa=0.4, metric=ICa_rmse} column=mean/sd raw={mean=1.2730634674060668, std=0.5597026136574934} n=5
\newcommand{\valGcaSweepICaRawPointFour}{1.273 \pm 0.560}

% source=results/gca_sweep_5seed/gca_sweep_multiseed.csv filter={gCa=0.4, metric=ICa_rmse} column=mean/sd raw={mean=1.2730634674060668, std=0.5597026136574934} n=5
\newcommand{\valGcaSweepICaRawPointFourMeanStd}{1.273 \pm 0.560}

% source=results/gca_sweep_5seed/gca_sweep_multiseed.csv filter={gCa=0.4, metric=ICa_rmse} column=mean raw=1.2730634674060668 n=5
\newcommand{\valGcaSweepICaRawPointFourMean}{1.273}

% source=results/gca_sweep_5seed/gca_sweep_multiseed.csv filter={gCa=0.4, metric=ICa_rmse} column=sd raw=0.5597026136574934 n=5
\newcommand{\valGcaSweepICaRawPointFourStd}{0.560}

% source=results/gca_sweep_5seed/gca_sweep_multiseed.csv filter={gCa=0.4, metric=ICa_rmse} column=n raw=5 n=5
\newcommand{\valGcaSweepICaRawPointFourN}{5}

% note: Voltage is forecast equally well at every conductance (0.19 to 0.25 across the sweep) while calcium recovery varies by 8x. Predictive accuracy on the observed channel is independent of whether the hidden term is identified, now backed by 5 seeds at every sweep point rather than one.
% source=results/gca_sweep_5seed/gca_sweep_multiseed.csv filter={gCa=0.4, metric=V_rmse} column=mean/sd raw={mean=0.21536875059604435, std=0.15566334732544684} n=5
\newcommand{\valGcaSweepVRmsePointFour}{0.215 \pm 0.156}

% note: Voltage is forecast equally well at every conductance (0.19 to 0.25 across the sweep) while calcium recovery varies by 8x. Predictive accuracy on the observed channel is independent of whether the hidden term is identified, now backed by 5 seeds at every sweep point rather than one.
% source=results/gca_sweep_5seed/gca_sweep_multiseed.csv filter={gCa=0.4, metric=V_rmse} column=mean/sd raw={mean=0.21536875059604435, std=0.15566334732544684} n=5
\newcommand{\valGcaSweepVRmsePointFourMeanStd}{0.215 \pm 0.156}

% note: Voltage is forecast equally well at every conductance (0.19 to 0.25 across the sweep) while calcium recovery varies by 8x. Predictive accuracy on the observed channel is independent of whether the hidden term is identified, now backed by 5 seeds at every sweep point rather than one.
% source=results/gca_sweep_5seed/gca_sweep_multiseed.csv filter={gCa=0.4, metric=V_rmse} column=mean raw=0.21536875059604435 n=5
\newcommand{\valGcaSweepVRmsePointFourMean}{0.215}

% note: Voltage is forecast equally well at every conductance (0.19 to 0.25 across the sweep) while calcium recovery varies by 8x. Predictive accuracy on the observed channel is independent of whether the hidden term is identified, now backed by 5 seeds at every sweep point rather than one.
% source=results/gca_sweep_5seed/gca_sweep_multiseed.csv filter={gCa=0.4, metric=V_rmse} column=sd raw=0.15566334732544684 n=5
\newcommand{\valGcaSweepVRmsePointFourStd}{0.156}

% note: Voltage is forecast equally well at every conductance (0.19 to 0.25 across the sweep) while calcium recovery varies by 8x. Predictive accuracy on the observed channel is independent of whether the hidden term is identified, now backed by 5 seeds at every sweep point rather than one.
% source=results/gca_sweep_5seed/gca_sweep_multiseed.csv filter={gCa=0.4, metric=V_rmse} column=n raw=5 n=5
\newcommand{\valGcaSweepVRmsePointFourN}{5}

% note: *** QUOTE THE MEDIAN, NOT THE MEAN. *** The mean (73.2 +/- 124) is dominated by seed 4444, the single seed that still blows up. Four of five seeds land between 8.7 and 34.5 mV with exactly correct spike counts. The honest summary is "usually finite, occasionally catastrophic", and the median plus the per-seed range says that; the mean hides it.
% source=results/node_parity/node_parity_summary.csv filter={key=parity_auto, metric=V_rmse} column=median raw=16.875472200783065 n=1
\newcommand{\valNodeParityAutoVRmseMedian}{16.88}

% note: Against the published 0.55 +/- 0.65 and the parity-with-time 0.43 +/- 0.45. The UDE saturates the 70 ms cap on every seed with zero variance, so the remaining contrast is as much about RELIABILITY as about magnitude.
% source=results/node_parity/node_parity_summary.csv filter={key=parity_auto, metric=rollout_horizon_ms} column=mean/sd raw={mean=19.530332681017615, std=17.752973380473332} n=5
\newcommand{\valNodeParityAutoRollout}{19.53 \pm 17.75}

% note: Against the published 0.55 +/- 0.65 and the parity-with-time 0.43 +/- 0.45. The UDE saturates the 70 ms cap on every seed with zero variance, so the remaining contrast is as much about RELIABILITY as about magnitude.
% source=results/node_parity/node_parity_summary.csv filter={key=parity_auto, metric=rollout_horizon_ms} column=mean/sd raw={mean=19.530332681017615, std=17.752973380473332} n=5
\newcommand{\valNodeParityAutoRolloutMeanStd}{19.53 \pm 17.75}

% note: Against the published 0.55 +/- 0.65 and the parity-with-time 0.43 +/- 0.45. The UDE saturates the 70 ms cap on every seed with zero variance, so the remaining contrast is as much about RELIABILITY as about magnitude.
% source=results/node_parity/node_parity_summary.csv filter={key=parity_auto, metric=rollout_horizon_ms} column=mean raw=19.530332681017615 n=5
\newcommand{\valNodeParityAutoRolloutMean}{19.53}

% note: Against the published 0.55 +/- 0.65 and the parity-with-time 0.43 +/- 0.45. The UDE saturates the 70 ms cap on every seed with zero variance, so the remaining contrast is as much about RELIABILITY as about magnitude.
% source=results/node_parity/node_parity_summary.csv filter={key=parity_auto, metric=rollout_horizon_ms} column=sd raw=17.752973380473332 n=5
\newcommand{\valNodeParityAutoRolloutStd}{17.75}

% note: Against the published 0.55 +/- 0.65 and the parity-with-time 0.43 +/- 0.45. The UDE saturates the 70 ms cap on every seed with zero variance, so the remaining contrast is as much about RELIABILITY as about magnitude.
% source=results/node_parity/node_parity_summary.csv filter={key=parity_auto, metric=rollout_horizon_ms} column=n raw=5 n=5
\newcommand{\valNodeParityAutoRolloutN}{5}

% note: The budget control. Same architecture as published, the UDE's own budget, and the horizon does not improve: it drops slightly (0.548 to 0.431). This is what rules out "the baseline was under-trained".
% source=results/node_parity/node_parity_summary.csv filter={key=parity_time, metric=rollout_horizon_ms} column=mean/sd raw={mean=0.43052837573385505, std=0.44625261256326343} n=5
\newcommand{\valNodeParityTimeRollout}{0.431 \pm 0.446}

% note: The budget control. Same architecture as published, the UDE's own budget, and the horizon does not improve: it drops slightly (0.548 to 0.431). This is what rules out "the baseline was under-trained".
% source=results/node_parity/node_parity_summary.csv filter={key=parity_time, metric=rollout_horizon_ms} column=mean/sd raw={mean=0.43052837573385505, std=0.44625261256326343} n=5
\newcommand{\valNodeParityTimeRolloutMeanStd}{0.431 \pm 0.446}

% note: The budget control. Same architecture as published, the UDE's own budget, and the horizon does not improve: it drops slightly (0.548 to 0.431). This is what rules out "the baseline was under-trained".
% source=results/node_parity/node_parity_summary.csv filter={key=parity_time, metric=rollout_horizon_ms} column=mean raw=0.43052837573385505 n=5
\newcommand{\valNodeParityTimeRolloutMean}{0.431}

% note: The budget control. Same architecture as published, the UDE's own budget, and the horizon does not improve: it drops slightly (0.548 to 0.431). This is what rules out "the baseline was under-trained".
% source=results/node_parity/node_parity_summary.csv filter={key=parity_time, metric=rollout_horizon_ms} column=sd raw=0.44625261256326343 n=5
\newcommand{\valNodeParityTimeRolloutStd}{0.446}

% note: The budget control. Same architecture as published, the UDE's own budget, and the horizon does not improve: it drops slightly (0.548 to 0.431). This is what rules out "the baseline was under-trained".
% source=results/node_parity/node_parity_summary.csv filter={key=parity_time, metric=rollout_horizon_ms} column=n raw=5 n=5
\newcommand{\valNodeParityTimeRolloutN}{5}

% note: *** GUARDED - DO NOT PRINT THIS MACRO. *** Guard G2 (node-no-meanstd) forbids a mean +/- SD summary of a BLOWN-UP integration, and this is one: per-seed values run 275 to 11762 mV (42.8x), the sd is 59% of the mean, and rmse_safe() computed each one over an unreported finite subset of timesteps. A dispersion statistic on a diverged solve is not an error bar. Print NodeParityTimeVRmseMin/Max instead. The entry is kept rather than deleted so this reasoning stays attached to the quantity. The DIRECTION is still the finding: more optimisation of the time-indexed architecture made the forecast worse, not better.
% source=results/node_parity/node_parity_summary.csv filter={key=parity_time, metric=V_rmse} column=mean/sd raw={mean=7118.513854016398, std=4205.73115496801} n=5
\newcommand{\valNodeParityTimeVRmse}{7119 \pm 4206}

% note: *** GUARDED - DO NOT PRINT THIS MACRO. *** Guard G2 (node-no-meanstd) forbids a mean +/- SD summary of a BLOWN-UP integration, and this is one: per-seed values run 275 to 11762 mV (42.8x), the sd is 59% of the mean, and rmse_safe() computed each one over an unreported finite subset of timesteps. A dispersion statistic on a diverged solve is not an error bar. Print NodeParityTimeVRmseMin/Max instead. The entry is kept rather than deleted so this reasoning stays attached to the quantity. The DIRECTION is still the finding: more optimisation of the time-indexed architecture made the forecast worse, not better.
% source=results/node_parity/node_parity_summary.csv filter={key=parity_time, metric=V_rmse} column=mean/sd raw={mean=7118.513854016398, std=4205.73115496801} n=5
\newcommand{\valNodeParityTimeVRmseMeanStd}{7119 \pm 4206}

% note: *** GUARDED - DO NOT PRINT THIS MACRO. *** Guard G2 (node-no-meanstd) forbids a mean +/- SD summary of a BLOWN-UP integration, and this is one: per-seed values run 275 to 11762 mV (42.8x), the sd is 59% of the mean, and rmse_safe() computed each one over an unreported finite subset of timesteps. A dispersion statistic on a diverged solve is not an error bar. Print NodeParityTimeVRmseMin/Max instead. The entry is kept rather than deleted so this reasoning stays attached to the quantity. The DIRECTION is still the finding: more optimisation of the time-indexed architecture made the forecast worse, not better.
% source=results/node_parity/node_parity_summary.csv filter={key=parity_time, metric=V_rmse} column=mean raw=7118.513854016398 n=5
\newcommand{\valNodeParityTimeVRmseMean}{7119}

% note: *** GUARDED - DO NOT PRINT THIS MACRO. *** Guard G2 (node-no-meanstd) forbids a mean +/- SD summary of a BLOWN-UP integration, and this is one: per-seed values run 275 to 11762 mV (42.8x), the sd is 59% of the mean, and rmse_safe() computed each one over an unreported finite subset of timesteps. A dispersion statistic on a diverged solve is not an error bar. Print NodeParityTimeVRmseMin/Max instead. The entry is kept rather than deleted so this reasoning stays attached to the quantity. The DIRECTION is still the finding: more optimisation of the time-indexed architecture made the forecast worse, not better.
% source=results/node_parity/node_parity_summary.csv filter={key=parity_time, metric=V_rmse} column=sd raw=4205.73115496801 n=5
\newcommand{\valNodeParityTimeVRmseStd}{4206}

% note: *** GUARDED - DO NOT PRINT THIS MACRO. *** Guard G2 (node-no-meanstd) forbids a mean +/- SD summary of a BLOWN-UP integration, and this is one: per-seed values run 275 to 11762 mV (42.8x), the sd is 59% of the mean, and rmse_safe() computed each one over an unreported finite subset of timesteps. A dispersion statistic on a diverged solve is not an error bar. Print NodeParityTimeVRmseMin/Max instead. The entry is kept rather than deleted so this reasoning stays attached to the quantity. The DIRECTION is still the finding: more optimisation of the time-indexed architecture made the forecast worse, not better.
% source=results/node_parity/node_parity_summary.csv filter={key=parity_time, metric=V_rmse} column=n raw=5 n=5
\newcommand{\valNodeParityTimeVRmseN}{5}

% source=results/node_parity/node_parity_summary.csv filter={key=parity_time, metric=V_rmse} column=min raw=274.5176471916351 n=1
\newcommand{\valNodeParityTimeVRmseMin}{275}

% note: Pair with NodeParityTimeVRmseMin. This is the admissible way to report the quantity: a span, which tells the reader the integration failed differently on every seed rather than by a characteristic amount.
% source=results/node_parity/node_parity_summary.csv filter={key=parity_time, metric=V_rmse} column=max raw=11761.564423376221 n=1
\newcommand{\valNodeParityTimeVRmseMax}{11762}

% source=results/identifiability/gca04_retrain_control.csv filter={case=before, domain=hull} column=a_hat raw=-0.37197985748605883 n=1
\newcommand{\valGcaFourRetrainAHatBefore}{-0.372}

% note: The physiological failure does not improve under a fourfold budget; the conductance moves further from the truth and stays wrong-signed. Pair with GcaFourRetrainRTwoAfter, which collapses below zero.
% source=results/identifiability/gca04_retrain_control.csv filter={case=after, domain=hull} column=a_hat raw=-1.341351055239545 n=1
\newcommand{\valGcaFourRetrainAHatAfter}{-1.341}

% source=results/identifiability/gca04_retrain_control.csv filter={case=before, domain=hull} column=r2_vs_true_on_domain raw=0.6881648342515774 n=1
\newcommand{\valGcaFourRetrainRTwoBefore}{0.688}

% note: NEGATIVE R^2: after the aggressive retrain the distilled conductance form is a worse predictor of the true calcium current than that current's own mean. State it plainly; it is the sharpest single number in the negative control.
% source=results/identifiability/gca04_retrain_control.csv filter={case=after, domain=hull} column=r2_vs_true_on_domain raw=-0.33137544146641895 n=1
\newcommand{\valGcaFourRetrainRTwoAfter}{-0.331}

% source=results/identifiability/voltage_only_commoneval.csv filter={model=UDE, observed=full, window=common_eval, metric=V_rmse} column=value (aggregated over seed) raw={mean=0.2663377958273806, std=0.0369932726112455} n=5
\newcommand{\valCommonEvalVRmseFull}{0.266 \pm 0.037}

% source=results/identifiability/voltage_only_commoneval.csv filter={model=UDE, observed=full, window=common_eval, metric=V_rmse} column=value (aggregated over seed) raw={mean=0.2663377958273806, std=0.0369932726112455} n=5
\newcommand{\valCommonEvalVRmseFullMeanStd}{0.266 \pm 0.037}

% source=results/identifiability/voltage_only_commoneval.csv filter={model=UDE, observed=full, window=common_eval, metric=V_rmse} column=value raw=0.2663377958273806 n=5
\newcommand{\valCommonEvalVRmseFullMean}{0.266}

% source=results/identifiability/voltage_only_commoneval.csv filter={model=UDE, observed=full, window=common_eval, metric=V_rmse} column=value raw=0.0369932726112455 n=5
\newcommand{\valCommonEvalVRmseFullStd}{0.037}

% source=results/identifiability/voltage_only_commoneval.csv filter={model=UDE, observed=full, window=common_eval, metric=V_rmse} column=seed raw=5 n=5
\newcommand{\valCommonEvalVRmseFullN}{5}

% note: Pair with CommonEvalVRmseFull. The two overlap comfortably, so the voltage-only equivalence holds on an interval that is UNSEEN for both conditions, not merely on each model's own forecast window. This is what the tombstone asked for. Same caveat as everywhere else in this paper: n=5, no statistical test, an equivalence-within-noise statement and not an ordering - and note the voltage-only SD is the larger of the two.
% source=results/identifiability/voltage_only_commoneval.csv filter={model=UDE, observed=voltage, window=common_eval, metric=V_rmse} column=value (aggregated over seed) raw={mean=0.2893368057271903, std=0.13672054332978156} n=5
\newcommand{\valCommonEvalVRmseVoltage}{0.289 \pm 0.137}

% note: Pair with CommonEvalVRmseFull. The two overlap comfortably, so the voltage-only equivalence holds on an interval that is UNSEEN for both conditions, not merely on each model's own forecast window. This is what the tombstone asked for. Same caveat as everywhere else in this paper: n=5, no statistical test, an equivalence-within-noise statement and not an ordering - and note the voltage-only SD is the larger of the two.
% source=results/identifiability/voltage_only_commoneval.csv filter={model=UDE, observed=voltage, window=common_eval, metric=V_rmse} column=value (aggregated over seed) raw={mean=0.2893368057271903, std=0.13672054332978156} n=5
\newcommand{\valCommonEvalVRmseVoltageMeanStd}{0.289 \pm 0.137}

% note: Pair with CommonEvalVRmseFull. The two overlap comfortably, so the voltage-only equivalence holds on an interval that is UNSEEN for both conditions, not merely on each model's own forecast window. This is what the tombstone asked for. Same caveat as everywhere else in this paper: n=5, no statistical test, an equivalence-within-noise statement and not an ordering - and note the voltage-only SD is the larger of the two.
% source=results/identifiability/voltage_only_commoneval.csv filter={model=UDE, observed=voltage, window=common_eval, metric=V_rmse} column=value raw=0.2893368057271903 n=5
\newcommand{\valCommonEvalVRmseVoltageMean}{0.289}

% note: Pair with CommonEvalVRmseFull. The two overlap comfortably, so the voltage-only equivalence holds on an interval that is UNSEEN for both conditions, not merely on each model's own forecast window. This is what the tombstone asked for. Same caveat as everywhere else in this paper: n=5, no statistical test, an equivalence-within-noise statement and not an ordering - and note the voltage-only SD is the larger of the two.
% source=results/identifiability/voltage_only_commoneval.csv filter={model=UDE, observed=voltage, window=common_eval, metric=V_rmse} column=value raw=0.13672054332978156 n=5
\newcommand{\valCommonEvalVRmseVoltageStd}{0.137}

% note: Pair with CommonEvalVRmseFull. The two overlap comfortably, so the voltage-only equivalence holds on an interval that is UNSEEN for both conditions, not merely on each model's own forecast window. This is what the tombstone asked for. Same caveat as everywhere else in this paper: n=5, no statistical test, an equivalence-within-noise statement and not an ordering - and note the voltage-only SD is the larger of the two.
% source=results/identifiability/voltage_only_commoneval.csv filter={model=UDE, observed=voltage, window=common_eval, metric=V_rmse} column=seed raw=5 n=5
\newcommand{\valCommonEvalVRmseVoltageN}{5}

% source=results/identifiability/voltage_only_commoneval.csv filter={model=UDE, observed=full, window=common_eval, metric=ICa_rmse} column=value (aggregated over seed) raw={mean=1.250596957705621, std=0.21188066901404193} n=5
\newcommand{\valCommonEvalICaRmseFull}{1.251 \pm 0.212}

% source=results/identifiability/voltage_only_commoneval.csv filter={model=UDE, observed=full, window=common_eval, metric=ICa_rmse} column=value (aggregated over seed) raw={mean=1.250596957705621, std=0.21188066901404193} n=5
\newcommand{\valCommonEvalICaRmseFullMeanStd}{1.251 \pm 0.212}

% source=results/identifiability/voltage_only_commoneval.csv filter={model=UDE, observed=full, window=common_eval, metric=ICa_rmse} column=value raw=1.250596957705621 n=5
\newcommand{\valCommonEvalICaRmseFullMean}{1.251}

% source=results/identifiability/voltage_only_commoneval.csv filter={model=UDE, observed=full, window=common_eval, metric=ICa_rmse} column=value raw=0.21188066901404193 n=5
\newcommand{\valCommonEvalICaRmseFullStd}{0.212}

% source=results/identifiability/voltage_only_commoneval.csv filter={model=UDE, observed=full, window=common_eval, metric=ICa_rmse} column=seed raw=5 n=5
\newcommand{\valCommonEvalICaRmseFullN}{5}

% source=results/identifiability/voltage_only_commoneval.csv filter={model=UDE, observed=voltage, window=common_eval, metric=ICa_rmse} column=value (aggregated over seed) raw={mean=1.3823637647825644, std=0.4777675534319743} n=5
\newcommand{\valCommonEvalICaRmseVoltage}{1.382 \pm 0.478}

% source=results/identifiability/voltage_only_commoneval.csv filter={model=UDE, observed=voltage, window=common_eval, metric=ICa_rmse} column=value (aggregated over seed) raw={mean=1.3823637647825644, std=0.4777675534319743} n=5
\newcommand{\valCommonEvalICaRmseVoltageMeanStd}{1.382 \pm 0.478}

% source=results/identifiability/voltage_only_commoneval.csv filter={model=UDE, observed=voltage, window=common_eval, metric=ICa_rmse} column=value raw=1.3823637647825644 n=5
\newcommand{\valCommonEvalICaRmseVoltageMean}{1.382}

% source=results/identifiability/voltage_only_commoneval.csv filter={model=UDE, observed=voltage, window=common_eval, metric=ICa_rmse} column=value raw=0.4777675534319743 n=5
\newcommand{\valCommonEvalICaRmseVoltageStd}{0.478}

% source=results/identifiability/voltage_only_commoneval.csv filter={model=UDE, observed=voltage, window=common_eval, metric=ICa_rmse} column=seed raw=5 n=5
\newcommand{\valCommonEvalICaRmseVoltageN}{5}

% source=results/node_parity/node_parity_summary.csv filter={key=published_time, metric=V_rmse} column=min raw=229.12221959973417 n=1
\newcommand{\valNodePublishedVRmseMin}{229}

% note: Same guard as NodeParityTimeVRmseMax. Supersedes NodeForecastVRmse's mean +/- SD wherever a magnitude is actually printed.
% source=results/node_parity/node_parity_summary.csv filter={key=published_time, metric=V_rmse} column=max raw=9213.415361612118 n=1
\newcommand{\valNodePublishedVRmseMax}{9213}
